\documentclass{article} % For LaTeX2e

% Lưu lại định nghĩa gốc của \addcontentsline trước khi style file vô hiệu hoá nó
\makeatletter
\let\orig@addcontentsline\addcontentsline
\makeatother

\usepackage{iclr2026_conference,times}

% Khôi phục lại \addcontentsline chuẩn cho mục lục
\makeatletter
\let\addcontentsline\orig@addcontentsline
\makeatother

\usepackage[utf8]{vietnam}
\usepackage{enumitem}
\usepackage{booktabs}
\usepackage{amsmath,amssymb}
\usepackage{graphicx}
\graphicspath{{images/}}
\usepackage{tikz}
\usetikzlibrary{calc}
\usepackage{xcolor}
\usepackage{tabularx}
\usepackage{array}
\usepackage{hyperref}
\usepackage{url}

% Optional math commands from https://github.com/goodfeli/dlbook_notation.
%%%%% NEW MATH DEFINITIONS %%%%%

\usepackage{amsmath,amsfonts,bm}

% Mark sections of captions for referring to divisions of figures
\newcommand{\figleft}{{\em (Left)}}
\newcommand{\figcenter}{{\em (Center)}}
\newcommand{\figright}{{\em (Right)}}
\newcommand{\figtop}{{\em (Top)}}
\newcommand{\figbottom}{{\em (Bottom)}}
\newcommand{\captiona}{{\em (a)}}
\newcommand{\captionb}{{\em (b)}}
\newcommand{\captionc}{{\em (c)}}
\newcommand{\captiond}{{\em (d)}}

% Highlight a newly defined term
\newcommand{\newterm}[1]{{\bf #1}}


% Figure reference, lower-case.
\def\figref#1{figure~\ref{#1}}
% Figure reference, capital. For start of sentence
\def\Figref#1{Figure~\ref{#1}}
\def\twofigref#1#2{figures \ref{#1} and \ref{#2}}
\def\quadfigref#1#2#3#4{figures \ref{#1}, \ref{#2}, \ref{#3} and \ref{#4}}
% Section reference, lower-case.
\def\secref#1{section~\ref{#1}}
% Section reference, capital.
\def\Secref#1{Section~\ref{#1}}
% Reference to two sections.
\def\twosecrefs#1#2{sections \ref{#1} and \ref{#2}}
% Reference to three sections.
\def\secrefs#1#2#3{sections \ref{#1}, \ref{#2} and \ref{#3}}
% Reference to an equation, lower-case.
\def\eqref#1{equation~\ref{#1}}
% Reference to an equation, upper case
\def\Eqref#1{Equation~\ref{#1}}
% A raw reference to an equation---avoid using if possible
\def\plaineqref#1{\ref{#1}}
% Reference to a chapter, lower-case.
\def\chapref#1{chapter~\ref{#1}}
% Reference to an equation, upper case.
\def\Chapref#1{Chapter~\ref{#1}}
% Reference to a range of chapters
\def\rangechapref#1#2{chapters\ref{#1}--\ref{#2}}
% Reference to an algorithm, lower-case.
\def\algref#1{algorithm~\ref{#1}}
% Reference to an algorithm, upper case.
\def\Algref#1{Algorithm~\ref{#1}}
\def\twoalgref#1#2{algorithms \ref{#1} and \ref{#2}}
\def\Twoalgref#1#2{Algorithms \ref{#1} and \ref{#2}}
% Reference to a part, lower case
\def\partref#1{part~\ref{#1}}
% Reference to a part, upper case
\def\Partref#1{Part~\ref{#1}}
\def\twopartref#1#2{parts \ref{#1} and \ref{#2}}

\def\ceil#1{\lceil #1 \rceil}
\def\floor#1{\lfloor #1 \rfloor}
\def\1{\bm{1}}
\newcommand{\train}{\mathcal{D}}
\newcommand{\valid}{\mathcal{D_{\mathrm{valid}}}}
\newcommand{\test}{\mathcal{D_{\mathrm{test}}}}

\def\eps{{\epsilon}}


% Random variables
\def\reta{{\textnormal{$\eta$}}}
\def\ra{{\textnormal{a}}}
\def\rb{{\textnormal{b}}}
\def\rc{{\textnormal{c}}}
\def\rd{{\textnormal{d}}}
\def\re{{\textnormal{e}}}
\def\rf{{\textnormal{f}}}
\def\rg{{\textnormal{g}}}
\def\rh{{\textnormal{h}}}
\def\ri{{\textnormal{i}}}
\def\rj{{\textnormal{j}}}
\def\rk{{\textnormal{k}}}
\def\rl{{\textnormal{l}}}
% rm is already a command, just don't name any random variables m
\def\rn{{\textnormal{n}}}
\def\ro{{\textnormal{o}}}
\def\rp{{\textnormal{p}}}
\def\rq{{\textnormal{q}}}
\def\rr{{\textnormal{r}}}
\def\rs{{\textnormal{s}}}
\def\rt{{\textnormal{t}}}
\def\ru{{\textnormal{u}}}
\def\rv{{\textnormal{v}}}
\def\rw{{\textnormal{w}}}
\def\rx{{\textnormal{x}}}
\def\ry{{\textnormal{y}}}
\def\rz{{\textnormal{z}}}

% Random vectors
\def\rvepsilon{{\mathbf{\epsilon}}}
\def\rvtheta{{\mathbf{\theta}}}
\def\rva{{\mathbf{a}}}
\def\rvb{{\mathbf{b}}}
\def\rvc{{\mathbf{c}}}
\def\rvd{{\mathbf{d}}}
\def\rve{{\mathbf{e}}}
\def\rvf{{\mathbf{f}}}
\def\rvg{{\mathbf{g}}}
\def\rvh{{\mathbf{h}}}
\def\rvu{{\mathbf{i}}}
\def\rvj{{\mathbf{j}}}
\def\rvk{{\mathbf{k}}}
\def\rvl{{\mathbf{l}}}
\def\rvm{{\mathbf{m}}}
\def\rvn{{\mathbf{n}}}
\def\rvo{{\mathbf{o}}}
\def\rvp{{\mathbf{p}}}
\def\rvq{{\mathbf{q}}}
\def\rvr{{\mathbf{r}}}
\def\rvs{{\mathbf{s}}}
\def\rvt{{\mathbf{t}}}
\def\rvu{{\mathbf{u}}}
\def\rvv{{\mathbf{v}}}
\def\rvw{{\mathbf{w}}}
\def\rvx{{\mathbf{x}}}
\def\rvy{{\mathbf{y}}}
\def\rvz{{\mathbf{z}}}

% Elements of random vectors
\def\erva{{\textnormal{a}}}
\def\ervb{{\textnormal{b}}}
\def\ervc{{\textnormal{c}}}
\def\ervd{{\textnormal{d}}}
\def\erve{{\textnormal{e}}}
\def\ervf{{\textnormal{f}}}
\def\ervg{{\textnormal{g}}}
\def\ervh{{\textnormal{h}}}
\def\ervi{{\textnormal{i}}}
\def\ervj{{\textnormal{j}}}
\def\ervk{{\textnormal{k}}}
\def\ervl{{\textnormal{l}}}
\def\ervm{{\textnormal{m}}}
\def\ervn{{\textnormal{n}}}
\def\ervo{{\textnormal{o}}}
\def\ervp{{\textnormal{p}}}
\def\ervq{{\textnormal{q}}}
\def\ervr{{\textnormal{r}}}
\def\ervs{{\textnormal{s}}}
\def\ervt{{\textnormal{t}}}
\def\ervu{{\textnormal{u}}}
\def\ervv{{\textnormal{v}}}
\def\ervw{{\textnormal{w}}}
\def\ervx{{\textnormal{x}}}
\def\ervy{{\textnormal{y}}}
\def\ervz{{\textnormal{z}}}

% Random matrices
\def\rmA{{\mathbf{A}}}
\def\rmB{{\mathbf{B}}}
\def\rmC{{\mathbf{C}}}
\def\rmD{{\mathbf{D}}}
\def\rmE{{\mathbf{E}}}
\def\rmF{{\mathbf{F}}}
\def\rmG{{\mathbf{G}}}
\def\rmH{{\mathbf{H}}}
\def\rmI{{\mathbf{I}}}
\def\rmJ{{\mathbf{J}}}
\def\rmK{{\mathbf{K}}}
\def\rmL{{\mathbf{L}}}
\def\rmM{{\mathbf{M}}}
\def\rmN{{\mathbf{N}}}
\def\rmO{{\mathbf{O}}}
\def\rmP{{\mathbf{P}}}
\def\rmQ{{\mathbf{Q}}}
\def\rmR{{\mathbf{R}}}
\def\rmS{{\mathbf{S}}}
\def\rmT{{\mathbf{T}}}
\def\rmU{{\mathbf{U}}}
\def\rmV{{\mathbf{V}}}
\def\rmW{{\mathbf{W}}}
\def\rmX{{\mathbf{X}}}
\def\rmY{{\mathbf{Y}}}
\def\rmZ{{\mathbf{Z}}}

% Elements of random matrices
\def\ermA{{\textnormal{A}}}
\def\ermB{{\textnormal{B}}}
\def\ermC{{\textnormal{C}}}
\def\ermD{{\textnormal{D}}}
\def\ermE{{\textnormal{E}}}
\def\ermF{{\textnormal{F}}}
\def\ermG{{\textnormal{G}}}
\def\ermH{{\textnormal{H}}}
\def\ermI{{\textnormal{I}}}
\def\ermJ{{\textnormal{J}}}
\def\ermK{{\textnormal{K}}}
\def\ermL{{\textnormal{L}}}
\def\ermM{{\textnormal{M}}}
\def\ermN{{\textnormal{N}}}
\def\ermO{{\textnormal{O}}}
\def\ermP{{\textnormal{P}}}
\def\ermQ{{\textnormal{Q}}}
\def\ermR{{\textnormal{R}}}
\def\ermS{{\textnormal{S}}}
\def\ermT{{\textnormal{T}}}
\def\ermU{{\textnormal{U}}}
\def\ermV{{\textnormal{V}}}
\def\ermW{{\textnormal{W}}}
\def\ermX{{\textnormal{X}}}
\def\ermY{{\textnormal{Y}}}
\def\ermZ{{\textnormal{Z}}}

% Vectors
\def\vzero{{\bm{0}}}
\def\vone{{\bm{1}}}
\def\vmu{{\bm{\mu}}}
\def\vtheta{{\bm{\theta}}}
\def\va{{\bm{a}}}
\def\vb{{\bm{b}}}
\def\vc{{\bm{c}}}
\def\vd{{\bm{d}}}
\def\ve{{\bm{e}}}
\def\vf{{\bm{f}}}
\def\vg{{\bm{g}}}
\def\vh{{\bm{h}}}
\def\vi{{\bm{i}}}
\def\vj{{\bm{j}}}
\def\vk{{\bm{k}}}
\def\vl{{\bm{l}}}
\def\vm{{\bm{m}}}
\def\vn{{\bm{n}}}
\def\vo{{\bm{o}}}
\def\vp{{\bm{p}}}
\def\vq{{\bm{q}}}
\def\vr{{\bm{r}}}
\def\vs{{\bm{s}}}
\def\vt{{\bm{t}}}
\def\vu{{\bm{u}}}
\def\vv{{\bm{v}}}
\def\vw{{\bm{w}}}
\def\vx{{\bm{x}}}
\def\vy{{\bm{y}}}
\def\vz{{\bm{z}}}

% Elements of vectors
\def\evalpha{{\alpha}}
\def\evbeta{{\beta}}
\def\evepsilon{{\epsilon}}
\def\evlambda{{\lambda}}
\def\evomega{{\omega}}
\def\evmu{{\mu}}
\def\evpsi{{\psi}}
\def\evsigma{{\sigma}}
\def\evtheta{{\theta}}
\def\eva{{a}}
\def\evb{{b}}
\def\evc{{c}}
\def\evd{{d}}
\def\eve{{e}}
\def\evf{{f}}
\def\evg{{g}}
\def\evh{{h}}
\def\evi{{i}}
\def\evj{{j}}
\def\evk{{k}}
\def\evl{{l}}
\def\evm{{m}}
\def\evn{{n}}
\def\evo{{o}}
\def\evp{{p}}
\def\evq{{q}}
\def\evr{{r}}
\def\evs{{s}}
\def\evt{{t}}
\def\evu{{u}}
\def\evv{{v}}
\def\evw{{w}}
\def\evx{{x}}
\def\evy{{y}}
\def\evz{{z}}

% Matrix
\def\mA{{\bm{A}}}
\def\mB{{\bm{B}}}
\def\mC{{\bm{C}}}
\def\mD{{\bm{D}}}
\def\mE{{\bm{E}}}
\def\mF{{\bm{F}}}
\def\mG{{\bm{G}}}
\def\mH{{\bm{H}}}
\def\mI{{\bm{I}}}
\def\mJ{{\bm{J}}}
\def\mK{{\bm{K}}}
\def\mL{{\bm{L}}}
\def\mM{{\bm{M}}}
\def\mN{{\bm{N}}}
\def\mO{{\bm{O}}}
\def\mP{{\bm{P}}}
\def\mQ{{\bm{Q}}}
\def\mR{{\bm{R}}}
\def\mS{{\bm{S}}}
\def\mT{{\bm{T}}}
\def\mU{{\bm{U}}}
\def\mV{{\bm{V}}}
\def\mW{{\bm{W}}}
\def\mX{{\bm{X}}}
\def\mY{{\bm{Y}}}
\def\mZ{{\bm{Z}}}
\def\mBeta{{\bm{\beta}}}
\def\mPhi{{\bm{\Phi}}}
\def\mLambda{{\bm{\Lambda}}}
\def\mSigma{{\bm{\Sigma}}}

% Tensor
\DeclareMathAlphabet{\mathsfit}{\encodingdefault}{\sfdefault}{m}{sl}
\SetMathAlphabet{\mathsfit}{bold}{\encodingdefault}{\sfdefault}{bx}{n}
\newcommand{\tens}[1]{\bm{\mathsfit{#1}}}
\def\tA{{\tens{A}}}
\def\tB{{\tens{B}}}
\def\tC{{\tens{C}}}
\def\tD{{\tens{D}}}
\def\tE{{\tens{E}}}
\def\tF{{\tens{F}}}
\def\tG{{\tens{G}}}
\def\tH{{\tens{H}}}
\def\tI{{\tens{I}}}
\def\tJ{{\tens{J}}}
\def\tK{{\tens{K}}}
\def\tL{{\tens{L}}}
\def\tM{{\tens{M}}}
\def\tN{{\tens{N}}}
\def\tO{{\tens{O}}}
\def\tP{{\tens{P}}}
\def\tQ{{\tens{Q}}}
\def\tR{{\tens{R}}}
\def\tS{{\tens{S}}}
\def\tT{{\tens{T}}}
\def\tU{{\tens{U}}}
\def\tV{{\tens{V}}}
\def\tW{{\tens{W}}}
\def\tX{{\tens{X}}}
\def\tY{{\tens{Y}}}
\def\tZ{{\tens{Z}}}


% Graph
\def\gA{{\mathcal{A}}}
\def\gB{{\mathcal{B}}}
\def\gC{{\mathcal{C}}}
\def\gD{{\mathcal{D}}}
\def\gE{{\mathcal{E}}}
\def\gF{{\mathcal{F}}}
\def\gG{{\mathcal{G}}}
\def\gH{{\mathcal{H}}}
\def\gI{{\mathcal{I}}}
\def\gJ{{\mathcal{J}}}
\def\gK{{\mathcal{K}}}
\def\gL{{\mathcal{L}}}
\def\gM{{\mathcal{M}}}
\def\gN{{\mathcal{N}}}
\def\gO{{\mathcal{O}}}
\def\gP{{\mathcal{P}}}
\def\gQ{{\mathcal{Q}}}
\def\gR{{\mathcal{R}}}
\def\gS{{\mathcal{S}}}
\def\gT{{\mathcal{T}}}
\def\gU{{\mathcal{U}}}
\def\gV{{\mathcal{V}}}
\def\gW{{\mathcal{W}}}
\def\gX{{\mathcal{X}}}
\def\gY{{\mathcal{Y}}}
\def\gZ{{\mathcal{Z}}}

% Sets
\def\sA{{\mathbb{A}}}
\def\sB{{\mathbb{B}}}
\def\sC{{\mathbb{C}}}
\def\sD{{\mathbb{D}}}
% Don't use a set called E, because this would be the same as our symbol
% for expectation.
\def\sF{{\mathbb{F}}}
\def\sG{{\mathbb{G}}}
\def\sH{{\mathbb{H}}}
\def\sI{{\mathbb{I}}}
\def\sJ{{\mathbb{J}}}
\def\sK{{\mathbb{K}}}
\def\sL{{\mathbb{L}}}
\def\sM{{\mathbb{M}}}
\def\sN{{\mathbb{N}}}
\def\sO{{\mathbb{O}}}
\def\sP{{\mathbb{P}}}
\def\sQ{{\mathbb{Q}}}
\def\sR{{\mathbb{R}}}
\def\sS{{\mathbb{S}}}
\def\sT{{\mathbb{T}}}
\def\sU{{\mathbb{U}}}
\def\sV{{\mathbb{V}}}
\def\sW{{\mathbb{W}}}
\def\sX{{\mathbb{X}}}
\def\sY{{\mathbb{Y}}}
\def\sZ{{\mathbb{Z}}}

% Entries of a matrix
\def\emLambda{{\Lambda}}
\def\emA{{A}}
\def\emB{{B}}
\def\emC{{C}}
\def\emD{{D}}
\def\emE{{E}}
\def\emF{{F}}
\def\emG{{G}}
\def\emH{{H}}
\def\emI{{I}}
\def\emJ{{J}}
\def\emK{{K}}
\def\emL{{L}}
\def\emM{{M}}
\def\emN{{N}}
\def\emO{{O}}
\def\emP{{P}}
\def\emQ{{Q}}
\def\emR{{R}}
\def\emS{{S}}
\def\emT{{T}}
\def\emU{{U}}
\def\emV{{V}}
\def\emW{{W}}
\def\emX{{X}}
\def\emY{{Y}}
\def\emZ{{Z}}
\def\emSigma{{\Sigma}}

% entries of a tensor
% Same font as tensor, without \bm wrapper
\newcommand{\etens}[1]{\mathsfit{#1}}
\def\etLambda{{\etens{\Lambda}}}
\def\etA{{\etens{A}}}
\def\etB{{\etens{B}}}
\def\etC{{\etens{C}}}
\def\etD{{\etens{D}}}
\def\etE{{\etens{E}}}
\def\etF{{\etens{F}}}
\def\etG{{\etens{G}}}
\def\etH{{\etens{H}}}
\def\etI{{\etens{I}}}
\def\etJ{{\etens{J}}}
\def\etK{{\etens{K}}}
\def\etL{{\etens{L}}}
\def\etM{{\etens{M}}}
\def\etN{{\etens{N}}}
\def\etO{{\etens{O}}}
\def\etP{{\etens{P}}}
\def\etQ{{\etens{Q}}}
\def\etR{{\etens{R}}}
\def\etS{{\etens{S}}}
\def\etT{{\etens{T}}}
\def\etU{{\etens{U}}}
\def\etV{{\etens{V}}}
\def\etW{{\etens{W}}}
\def\etX{{\etens{X}}}
\def\etY{{\etens{Y}}}
\def\etZ{{\etens{Z}}}

% The true underlying data generating distribution
\newcommand{\pdata}{p_{\rm{data}}}
% The empirical distribution defined by the training set
\newcommand{\ptrain}{\hat{p}_{\rm{data}}}
\newcommand{\Ptrain}{\hat{P}_{\rm{data}}}
% The model distribution
\newcommand{\pmodel}{p_{\rm{model}}}
\newcommand{\Pmodel}{P_{\rm{model}}}
\newcommand{\ptildemodel}{\tilde{p}_{\rm{model}}}
% Stochastic autoencoder distributions
\newcommand{\pencode}{p_{\rm{encoder}}}
\newcommand{\pdecode}{p_{\rm{decoder}}}
\newcommand{\precons}{p_{\rm{reconstruct}}}

\newcommand{\laplace}{\mathrm{Laplace}} % Laplace distribution

\newcommand{\E}{\mathbb{E}}
\newcommand{\Ls}{\mathcal{L}}
\newcommand{\R}{\mathbb{R}}
\newcommand{\emp}{\tilde{p}}
\newcommand{\lr}{\alpha}
\newcommand{\reg}{\lambda}
\newcommand{\rect}{\mathrm{rectifier}}
\newcommand{\softmax}{\mathrm{softmax}}
\newcommand{\sigmoid}{\sigma}
\newcommand{\softplus}{\zeta}
\newcommand{\KL}{D_{\mathrm{KL}}}
\newcommand{\Var}{\mathrm{Var}}
\newcommand{\standarderror}{\mathrm{SE}}
\newcommand{\Cov}{\mathrm{Cov}}
% Wolfram Mathworld says $L^2$ is for function spaces and $\ell^2$ is for vectors
% But then they seem to use $L^2$ for vectors throughout the site, and so does
% wikipedia.
\newcommand{\normlzero}{L^0}
\newcommand{\normlone}{L^1}
\newcommand{\normltwo}{L^2}
\newcommand{\normlp}{L^p}
\newcommand{\normmax}{L^\infty}

\newcommand{\parents}{Pa} % See usage in notation.tex. Chosen to match Daphne's book.

\DeclareMathOperator*{\argmax}{arg\,max}
\DeclareMathOperator*{\argmin}{arg\,min}

\DeclareMathOperator{\sign}{sign}
\DeclareMathOperator{\Tr}{Tr}
\let\ab\allowbreak


\setcounter{tocdepth}{2}

\title{BÁO CÁO BÀI TẬP LỚP NHÀ: KHOA HỌC DỮ LIỆU THỊ GIÁC VÀ ĐỒ HỌA MÁY TÍNH}

\author{
  Tô Thành Đạt -- 24127153 \\
  Khoa Công nghệ Thông tin \\
  Trường ĐH Khoa học Tự nhiên \\
  \texttt{24127153@student.hcmus.edu.vn} \\
  \And
  Lương Hưng Phát -- 24127298 \\
  Khoa Công nghệ Thông tin \\
  Trường ĐH Khoa học Tự nhiên \\
  \texttt{24127298@student.hcmus.edu.vn} \\
  \And
  Nguyễn Nhựt Huy -- 24127398 \\
  Khoa Công nghệ Thông tin \\
  Trường ĐH Khoa học Tự nhiên \\
  \texttt{24127398@student.hcmus.edu.vn} \\
}

\newcommand{\fix}{\marginpar{FIX}}
\newcommand{\new}{\marginpar{NEW}}

\iclrfinalcopy % Uncommented for full header display

\begin{document}

% -------------------------------
% Trang bìa
% -------------------------------
\begin{titlepage}
\begin{tikzpicture}[remember picture,overlay,inner sep=0,outer sep=0]
     \draw[blue!70!black,line width=4pt] ([xshift=-1.5cm,yshift=-2cm]current page.north east) coordinate (A)--([xshift=1.5cm,yshift=-2cm]current page.north west) coordinate(B)--([xshift=1.5cm,yshift=2cm]current page.south west) coordinate (C)--([xshift=-1.5cm,yshift=2cm]current page.south east) coordinate(D)--cycle;

     \draw ([yshift=0.5cm,xshift=-0.5cm]A)-- ([yshift=0.5cm,xshift=0.5cm]B)--
     ([yshift=-0.5cm,xshift=0.5cm]B) --([yshift=-0.5cm,xshift=-0.5cm]B)--([yshift=0.5cm,xshift=-0.5cm]C)--([yshift=0.5cm,xshift=0.5cm]C)--([yshift=-0.5cm,xshift=0.5cm]C)-- ([yshift=-0.5cm,xshift=-0.5cm]D)--([yshift=0.5cm,xshift=-0.5cm]D)--([yshift=0.5cm,xshift=0.5cm]D)--([yshift=-0.5cm,xshift=0.5cm]A)--([yshift=-0.5cm,xshift=-0.5cm]A)--([yshift=0.5cm,xshift=-0.5cm]A);


     \draw ([yshift=-0.3cm,xshift=0.3cm]A)-- ([yshift=-0.3cm,xshift=-0.3cm]B)--
     ([yshift=0.3cm,xshift=-0.3cm]B) --([yshift=0.3cm,xshift=0.3cm]B)--([yshift=-0.3cm,xshift=0.3cm]C)--([yshift=-0.3cm,xshift=-0.3cm]C)--([yshift=0.3cm,xshift=-0.3cm]C)-- ([yshift=0.3cm,xshift=0.3cm]D)--([yshift=-0.3cm,xshift=0.3cm]D)--([yshift=-0.3cm,xshift=-0.3cm]D)--([yshift=0.3cm,xshift=-0.3cm]A)--([yshift=0.3cm,xshift=0.3cm]A)--([yshift=-0.3cm,xshift=0.3cm]A);

   \end{tikzpicture}


  % HEADER
\begin{center}
    \normalfont\bfseries TRƯỜNG ĐẠI HỌC KHOA HỌC TỰ NHIÊN \\[0.1cm]
    \textbf{KHOA CÔNG NGHỆ THÔNG TIN}\\
    \includegraphics[scale=0.5]{hcmus_logo.png}\\[0.5cm]
    {\LARGE \textbf{ĐỒ HỌA MÁY TÍNH}} \\[0.3cm]
    {\Large \textbf{BÀI TẬP LỚP NHÀ}}
\end{center}


\vspace*{1.2cm}
 % MENTOR LIST
\hspace*{2cm}
{\small % hoặc \footnotesize / \scriptsize
\begin{tabular}{@{}l l@{}}
{\normalsize \textbf{Giảng viên:}} & Thầy Lý Quốc Ngọc
\end{tabular}
}

\vspace*{0.6cm}
  % AUTHOR LIST
\hspace*{2cm} {\normalsize \textbf{Sinh viên thực hiện:}}
\begin{center}
  \setlength{\tabcolsep}{1.8em}
  {\normalfont
    \renewcommand{\arraystretch}{0.9} % giảm khoảng cách dòng
    \small % hoặc footnotesize / scriptsize
    \begin{tabular}{@{}c c c@{}} % <-- 3 cột
      \textbf{Họ và tên} & \textbf{Mã số sinh viên} & \textbf{Lớp} \\
            Tô Thành Đạt    & 24127153 & 24TGMT \\
            Lương Hưng Phát & 24127298 & 24TGMT \\
            Nguyễn Nhựt Huy & 24127398 & 24TGMT
    \end{tabular}
  }
\end{center}

  % FOOTER DATE
\vspace*{2.5cm}
  \begin{center}
    TP. Hồ Chí Minh, tháng 10/2026
  \end{center}
\end{titlepage}



\clearpage

\maketitle

\begin{abstract}
Báo cáo này tổng hợp toàn bộ nội dung bài học, câu hỏi lý thuyết và bài tập môn Khoa học Dữ liệu Thị giác (Visual Data Science) và Đồ họa Máy tính (Computer Graphics) theo lộ trình 12 tuần học. Nội dung bao gồm các nền tảng về cuộc Cách mạng Công nghiệp 4.0, xử lý ảnh, thị giác máy tính cổ điển, các kiến trúc học sâu tiên tiến (CNNs, Transformers, Generative Models, Multimodal AI) cũng như các ứng dụng thực tiễn trong sản xuất và đời sống.
\end{abstract}

\tableofcontents
\newpage

% --- 12 Tuần học ---
\section{Tuần 1: Tổng quan Khoa học Dữ liệu Thị giác \& Computer Graphics}
\label{sec:week1}

\subsection{Đối tượng nghiên cứu của Khoa học Dữ liệu Thị giác và Dữ liệu Đa phương thức}

Khoa học Dữ liệu Thị giác (Visual Data Science) là lĩnh vực liên ngành nghiên cứu các nguyên lý toán học, cấu trúc dữ liệu và mô hình thuật toán nhằm thu thập, biểu diễn, trích xuất đặc trưng và suy luận thông tin từ dữ liệu trực quan. Dữ liệu này bao gồm hình ảnh số 2D, chuỗi khung hình video có tính liên tục thời gian, đám mây điểm 3D (Point Clouds) từ LiDAR, thể tích voxel trong chẩn đoán y khoa (CT, MRI), dữ liệu ảnh viễn thám đa phổ, v.v.

Bên cạnh dữ liệu thị giác thuần túy, ngành còn mở rộng nghiên cứu sang dữ liệu đa phương thức (Multimodal Data). Việc liên kết dữ liệu thị giác với các phương thức khác như văn bản tự nhiên, tín hiệu âm thanh hay chuỗi thời gian từ cảm biến quán tính (IMU) giúp giải quyết các bài toán biểu diễn ngữ cảnh phức tạp. Một ví dụ tiêu biểu là bài toán phân đoạn đối tượng trong video dựa trên mô tả hành động bằng văn bản hoặc âm thanh (Text/Audio-based Referring Motion Expression Video Segmentation)~\citep{week1_MeViS}. Thay vì chỉ định mục tiêu bằng các index của class, mô hình tiếp nhận câu mô tả tự nhiên hoặc âm thanh đặc trưng để khoanh vùng chính xác mặt nạ pixel của đối tượng qua từng khung hình.

\subsection{Mục tiêu và bản chất của cuộc Cách mạng Công nghiệp 4.0}

Trọng tâm của Cách mạng Công nghiệp 4.0 (CMCN 4.0) là tự động hóa thông minh (Smartification) quy trình sản xuất và vận hành xã hội trên quy mô lớn, dựa trên nền tảng chuyển đổi số.

Quá trình thông minh hóa biểu hiện ở khả năng tích hợp thuật toán học máy và mạng lưới cảm biến IoT vào thiết bị phần cứng, giúp hệ thống tự thu thập dữ liệu, tự hiệu chỉnh thông số vận hành và ra quyết định tối ưu theo thời gian thực mà không phụ thuộc vào sự can thiệp thủ công liên tục của con người.

Trong \textbf{sản xuất} công nghiệp, giải pháp này giúp nâng cao năng suất, giảm thiểu chi phí vận hành và loại bỏ sai sót chủ quan, đồng thời cho phép tùy biến sản phẩm hàng loạt (Mass Customization). Trong \textbf{quản lý xã hội}, các hệ thống điều khiển giao thông thích ứng, lưới điện thông minh và y tế số giúp nâng cao chất lượng dịch vụ công và phân bổ nguồn lực đô thị minh bạch.

Cần phân biệt rõ ba cấp độ chuyển đổi:
\begin{itemize}
    \item \textit{Số hóa dữ liệu (Digitization):} Chuyển đổi dữ liệu từ dạng vật lý hoặc tương tự (analog, giấy tờ, băng từ) sang định dạng nhị phân số (bit 0 và 1).
    \item \textit{Số hóa quy trình (Digitalization):} Ứng dụng công nghệ kỹ thuật số và dữ liệu đã số hóa vào việc cải tiến hoặc tự động hóa các quy trình nghiệp vụ sẵn có mà không làm thay đổi bản chất mô hình vận hành.
    \item \textit{Chuyển đổi số (Digital Transformation):} Tái cấu trúc toàn diện mô hình hoạt động, quy trình làm việc, chiến lược ra quyết định và văn hóa tổ chức dựa trên việc khai thác dữ liệu và công nghệ số thông minh.
\end{itemize}

Trong hệ sinh thái công nghệ 4.0, các ngành khoa học cơ bản như Toán học, Vật lý, Hóa học và Sinh học cung cấp cơ sở lý thuyết và chế tạo linh kiện bán dẫn, cảm biến quang học. Tuy nhiên, Trí tuệ nhân tạo (AI) giữ vai trò trung tâm vì là công cụ duy nhất có khả năng trích xuất tri thức từ các tập gồm lượng cực lớn dữ liệu phi cấu trúc và đưa ra quyết định trong môi trường biến động liên tục.

\subsection{Tiến trình phát triển kỹ thuật qua các cuộc Cách mạng Công nghiệp}

Lịch sử phát triển sản xuất của loài người gắn liền với 4 bước nhảy vọt về công nghệ:

\textbf{CMCN 1.0 (Cuối thế kỷ 18 - đầu thế kỷ 19):} Bắt đầu khi động cơ hơi nước được thương mại hóa và ứng dụng vào ngành dệt, thay thế sức kéo cơ bắp và động vật bằng cơ giới hóa, thúc đẩy hình thành các đô thị công nghiệp đầu tiên.

\textbf{CMCN 2.0 (Cuối thế kỷ 19 - đầu thế kỷ 20):} Đánh dấu sự ra đời của năng lượng điện, động cơ đốt trong và dây chuyền lắp ráp liên tục (Assembly Line), đưa sản xuất hàng loạt với giá thành thấp trở thành chuẩn mực công nghiệp, đồng thời phát triển mạng lưới viễn thông và giao thông hiện đại.

\textbf{CMCN 3.0 (Thập niên 1960 - cuối thế kỷ 20):} Chuyển trọng tâm sang chất bán dẫn, máy vi tính và mạng Internet. Các bộ điều khiển logic lập trình được (PLC) trong tự động hóa công nghiệp đã thay đổi căn bản cách thức vận hành nhà máy và mở ra kỷ nguyên số hóa thông tin.

\textbf{CMCN 4.0 (Từ đầu thế kỷ 21 đến nay):} Xây dựng các hệ thống không gian mạng - vật lý (Cyber-Physical Systems - CPS), nơi thế giới thực và mô hình số đồng bộ hai chiều thông qua AI, IoT, dữ liệu lớn và điện toán đám mây, hướng tới các nhà máy tự hành và đô thị thông minh.

\subsection{Vai trò của dữ liệu thị giác trong sản xuất và quản lý xã hội}

Thị giác cung cấp hơn 80\% lượng thông tin nhận thức của con người về thế giới xung quanh. Trong kỹ thuật, dữ liệu thị giác là phương thức thu thập thông tin không tiếp xúc (non-contact) và không phá hủy mẫu (non-destructive).

Trong \textbf{sản xuất}, hệ thống camera công nghiệp thu thập dữ liệu hình ảnh với tần số quét cao để giám sát độ phẳng, kích thước hình học và phát hiện khuyết tật vật liệu theo thời gian thực. Trong quản trị xã hội, dữ liệu ảnh và video từ hệ thống camera giám sát đô thị và ảnh viễn thám cung cấp thông tin không gian chi tiết, hỗ trợ trích xuất luồng phương tiện, phát hiện ùn tắc và phục vụ công tác cứu hộ, an ninh trật tự.

\subsection{Ứng dụng của dữ liệu thị giác trong công nghiệp, nông nghiệp và lâm nghiệp}

\textbf{Trong công nghiệp sản xuất \& chế tạo:}
\begin{itemize}
    \item \textit{Kiểm tra khuyết tật tự động (AOI):} Phát hiện vết xước bề mặt linh kiện, lỗi hàn vi mạch, lệch mối nối trên dây chuyền tốc độ cao với độ phân giải micromet.
    \item \textit{Thị giác robot (Vision-Guided Robotics):} Định vị tọa độ 3D để cánh tay robot tự động gắp thả linh kiện, hàn chính xác và lắp ráp phụ tùng.
    \item \textit{An toàn lao động:} Giám sát trang bị bảo hộ cá nhân (PPE), phát hiện công nhân đi vào vùng nguy hiểm quanh máy móc công nghiệp nặng.
\end{itemize}

\textbf{Trong nông nghiệp công nghệ cao:}
\begin{itemize}
    \item \textit{Phân loại nông sản:} Tự động phân loại kích thước, độ chín, màu sắc và loại bỏ dị tật trái cây sau thu hoạch.
    \item \textit{Bảo vệ thực vật:} Phát hiện sớm sâu bệnh và nấm lá qua camera trên cánh đồng, giúp phun thuốc chính xác theo vùng nhiễm bệnh.
    \item \textit{Giám sát mùa màng:} Dùng ảnh quang học và đa phổ từ drone để ước lượng chỉ số sinh khối (NDVI), dự báo năng suất thu hoạch.
\end{itemize}

\textbf{Trong lâm nghiệp \& Môi trường:}
\begin{itemize}
    \item \textit{Giám sát diện tích rừng:} Phân tích ảnh vệ tinh đa thời gian để theo dõi độ che phủ rừng, phát hiện hành vi khai thác gỗ trái phép.
    \item \textit{Cảnh báo sớm cháy rừng:} Hệ thống camera nhiệt và camera quang học tầm xa phát hiện cột khói và điểm nhiệt bất thường.
\end{itemize}

\subsection{Dữ liệu thị giác có được từ đâu? Sự khác biệt giữa mắt người và thiết bị thu nhận hình ảnh?}

Dữ liệu thị giác được thu thập từ nhiều nguồn cảm biến: camera quang học dải khả kiến (RGB cảm biến CCD/CMOS), camera nhiệt hồng ngoại, thiết bị quét laser LiDAR, cảm biến thời gian bay (ToF), máy chụp cắt lớp y tế (X-ray, CT, MRI, siêu âm) và dữ liệu mô phỏng 3D từ các engine đồ họa.

Cơ chế thị giác người và thiết bị thu nhận ảnh kỹ thuật số có những điểm khác biệt căn bản:
\begin{itemize}
    \item \textit{Cơ chế cảm nhận:} Mắt người sử dụng hàng triệu tế bào thụ cảm quang (tế bào que cho vùng tối và tế bào nón cho màu sắc) phân bố không đều trên võng mạc, tập trung mật độ cao nhất tại hoàng điểm (fovea). Tín hiệu được xử lý cục bộ tại võng mạc trước khi được truyền qua dây thần kinh thị giác lên vỏ não. Ngược lại, cảm biến máy ảnh kỹ thuật số dùng ma trận các điểm ảnh với bộ lọc màu Bayer (RGB), biến đổi photon thành điện tích tuyến tính tạo thành ma trận số thô.
    \item \textit{Dải động (Dynamic Range):} Mắt người có khả năng thích ứng động học linh hoạt vượt quá 20 stop ($>100,000:1$), giúp nhìn rõ đồng thời cả chi tiết trong bóng râm lẫn nguồn sáng gắt. Cảm biến thông thường chỉ đạt dải động khoảng 10--14 stop, dễ dẫn đến hiện tượng cháy sáng hoặc mất chi tiết vùng tối do nhiễu.
    \item \textit{Dải phổ (Spectral Range):} Mắt người chỉ cảm nhận được bước sóng trong khoảng $380\text - 740\text{ nm}$ (ánh sáng khả kiến). Thiết bị nhân tạo có thể thu nhận từ tia X, tia tử ngoại, hồng ngoại đến sóng radar vô tuyến với dải bước số rộng hơn.
    \item \textit{Tính định lượng vs Định tính:} Thiết bị quang học đo lường chính xác cường độ ánh sáng bằng các giá trị số tuyệt đối, vận hành liên tục không mệt mỏi. Mắt người phụ thuộc vào cảm nhận tâm lý học thị giác, mang tính định tính và dễ bị ảnh hưởng và suy giảm độ tập trung do mệt mỏi.
\end{itemize}

\subsection{Ngoài đặc trưng thị giác và ngữ nghĩa thì còn đặc trưng nào khác không?}

Bên cạnh các đặc trưng thị giác mức thấp (màu sắc, kết cấu texture, hình dạng) và đặc trưng ngữ nghĩa mức cao (nhãn lớp đối tượng), biểu diễn dữ liệu thị giác hiện đại bao gồm:

\textbf{Đặc trưng không gian (Spatial Features):} Chiều sâu pixel, tọa độ không gian thực, vector pháp tuyến bề mặt, tư thế góc xoay và tịnh tiến 6 bậc tự do, mô hình đám mây điểm và lưới đa giác (3D Mesh).

\textbf{Đặc trưng thời gian (Temporal Features):} Luồng quang học (Optical Flow) mô tả vector chuyển động của từng pixel giữa hai khung hình liên tiếp ($u, v$), quỹ đạo di chuyển (motion trajectory), vận tốc và gia tốc góc của thực thể.

\textbf{Đặc trưng cấu trúc (Relational Features):} Đồ thị quan hệ cảnh (Scene Graphs) biểu diễn các nút là thực thể và các cạnh là quan hệ không gian hoặc tương tác người - vật (Human-Object Interaction - HOI), giúp mô hình nắm bắt toàn cảnh bức tranh thay vì từng đối tượng đơn lẻ.

\textbf{Đặc trưng thống kê (Statistical Features):} Hệ số phổ Fourier biểu diễn tần số không gian cao/thấp, biến đổi Wavelet trích xuất cấu trúc đa tỷ lệ, ma trận đồng xuất mức xám (GLCM) và ma trận hiệp phương sai đặc trưng.

\subsection{Khái niệm đối tượng trong ảnh: Phân biệt cấp độ lớp và cấp độ thực thể}

Trong thị giác máy tính, một đối tượng thị giác (Visual Object) là một vùng không gian (hoặc tập hợp các pixel liên thông) tương ứng với một thực thể vật lý xác định trong thế giới thực, có ranh giới rõ ràng với hậu cảnh.

Có hai cấp độ định danh đối tượng:
\begin{itemize}
    \item \textit{Đối tượng chung (General Object / Category-level):} Nhận diện một thực thể thuộc về một lớp hay danh mục khái niệm chung (ví dụ: xác định một vùng ảnh là ``nam giới'', ``xe hơi'', hoặc ``con chó'').
    \item \textit{Đối tượng riêng (Specific Object / Instance-level):} Định danh chính xác một thực thể duy nhất trong vũ trụ đối tượng (ví dụ: nhận diện một người cụ thể là ``bạn An'', chiếc xe mang biển số riêng ``30F-999.99'', hoặc tòa nhà cụ thể ``Landmark 81'').
\end{itemize}

\subsection{Bản chất và thách thức của khoảng cách ngữ nghĩa (Semantic Gap)}

Khoảng cách ngữ nghĩa (Semantic Gap) là sự chênh lệch căn bản giữa cách biểu diễn dữ liệu mức thấp của máy tính và cách hiểu ngữ nghĩa mức cao của con người.

Về phía máy tính, hình ảnh chỉ là một mảng đa chiều chứa các giá trị số thô. Khi góc chiếu sáng, góc quan sát hoặc vị trí camera thay đổi, toàn bộ giá trị số trong tensor bị biến đổi hoàn toàn dù đối tượng thực tế không hề thay đổi.

Ngược lại, hệ thống thị giác con người tự động trích xuất các cấu trúc bất biến, kết hợp với tri thức đã có về thế giới quan để nhận thức ngay lập tức các thực thể, bối cảnh không gian, cảm xúc và các mối quan hệ logic phức tạp. Thu hẹp khoảng cách này là mục tiêu nghiên cứu trung tâm của thị giác máy tính.

\subsection{Tiến bộ hiện nay trong việc thu hẹp khoảng cách ngữ nghĩa}

Quá trình giải quyết Semantic Gap trong thị giác máy tính đã trải qua ba giai đoạn phát triển chính:
\begin{itemize}
    \item \textit{Giai đoạn trích xuất đặc trưng thủ công (trước 2012):} Sử dụng các thuật toán toán học xác định như SIFT, SURF, HOG kết hợp với bộ phân loại SVM hoặc Random Forest. Phương pháp này chỉ mô tả được các mẫu cục bộ đơn giản và phụ thuộc nặng nề vào việc tinh chỉnh đặc trưng bằng tay (hand-crafted features), khiến Semantic Gap vẫn rất lớn.
    \item \textit{Kỷ nguyên Học sâu với mạng CNN (2012--2020):} Sự xuất hiện của AlexNet, VGG và ResNet cho phép mạng tự động học biểu diễn đặc trưng phân tần. Điều này thu hẹp đáng kể khoảng cách ngữ nghĩa trong các tác vụ phân loại và phát hiện cơ bản.
    \item \textit{Kỷ nguyên Mô hình Nền tảng và Thị giác - Ngôn ngữ (2020 đến nay):} Kiến trúc Vision Transformer (ViT) kết hợp với các mô hình học biểu diễn liên kết đa phương thức (CLIP, BLIP, LLaVA, GPT-4V) đã ánh xạ trực tiếp không gian ảnh và không gian ngôn ngữ tự nhiên vào cùng một không gian nhúng chung (shared embedding space). Từ đó mà mô hình có thể suy luận ngữ cảnh, sinh mô tả tự nhiên (Image Captioning) và trả lời câu hỏi trực quan (Visual Question Answering).
\end{itemize}

Mặc dù đạt độ chính xác cao trên các tập dữ liệu chuẩn, các mô hình học sâu hiện nay vẫn rất nhạy cảm với nhiễu pixel đối kháng (adversarial perturbations), dễ sinh ảo giác (hallucination) trong điều kiện thiếu sáng hoặc bị che khuất một phần (partial occlusion), đồng thời còn hạn chế trong suy luận nhân quả (Causal Reasoning).

\subsection{Vì sao trong KHDL thị giác, ngoài dữ liệu hình ảnh và video thì ta còn phải mở rộng ra dữ liệu đa phương thức?}

Việc giới hạn mô hình trong dữ liệu hình ảnh hoặc video đơn lẻ (Unimodal) bộc lộ nhiều điểm nghẽn: Hình ảnh 2D thực chất là hình chiếu thu nhỏ của thế giới thực 3D, làm mất thông tin chiều sâu và dẫn đến tính nhập nhằng (ambiguity) về kích thước vật thể. Hơn nữa, nhiều hiện tượng vật lý trong thế giới thực đòi hỏi tín hiệu bổ trợ như âm thanh (ví dụ: tiếng va chạm xe hơi, tiếng rò rỉ khí), văn bản ngữ cảnh (biển báo, nhãn kỹ thuật) hay dữ liệu cảm biến đo đạc (nhiệt độ, gia tốc).

Dữ liệu đa phương thức cung cấp thêm thông tin cho dữ liệu thị giác và ngược lại. Ví dụ ngôn ngữ tự nhiên đóng vai trò là cầu nối ngữ nghĩa trừu tượng, trong khi dữ liệu thị giác cung cấp chi tiết hình học và vị trí không gian cụ thể. Sự kết hợp này là yêu cầu bắt buộc đối với các hệ thống tự hành như xe tự lái (phối hợp Camera RGB, Radar, LiDAR và bản đồ số HD) hay robot thông minh tương tác người - máy.

\subsection{Phân định ranh giới giữa Đồ họa máy tính, Xử lý ảnh và Thị giác máy tính}

Ba lĩnh vực này cùng thao tác trên dữ liệu thị giác nhưng có hướng chuyển đổi dữ liệu và mục tiêu hoàn toàn khác biệt. Bảng~\ref{tab:comparison_cg_ip_cv} so sánh tổng hợp giữa ba lĩnh vực trên các phương diện mục tiêu, dữ liệu đầu vào, dữ liệu đầu ra và bản chất chuyển đổi dữ liệu.

\begin{table}[ht]
\centering
\small
\caption{So sánh toàn diện giữa Computer Graphics, Image Processing và Computer Vision}
\label{tab:comparison_cg_ip_cv}
\renewcommand{\arraystretch}{1.25}
\begin{tabularx}{\textwidth}{|p{2.0cm}|X|X|X|}
\hline
\textbf{Tiêu chí} & \textbf{Computer Graphics (CG)} & \textbf{Image Processing (IP)} & \textbf{Computer Vision (CV)} \\
\hline
\textbf{Mục tiêu chính (Goal)} & Mô phỏng, kiến tạo và kết xuất thế giới thực hoặc ảo trực quan lên màn hình. & Nâng cao chất lượng, khôi phục, lọc nhiễu hoặc nén/biến đổi dữ liệu ảnh. & Hiểu, phân tích và trích xuất ngữ nghĩa, tri thức và đưa ra quyết định từ ảnh/video. \\
\hline
\textbf{Dữ liệu đầu vào (Input)} & Mô hình hình học 2D/3D (tập điểm, mesh, spline), tên thực thể, thuộc tính chất liệu, nguồn sáng. & Dữ liệu ảnh thô (\textit{Raw Image Data}), ảnh có nhiễu hoặc ảnh cần biến đổi. & Hình ảnh hoặc chuỗi video (\textit{Visual Data}) từ thế giới thực. \\
\hline
\textbf{Dữ liệu đầu ra (Output)} & Hình ảnh hoặc video được kết xuất trực quan (\textit{Rendering Output}). & Hình ảnh đã qua xử lý, tăng cường hoặc nén (\textit{Processed Image Data}). & Thông tin ngữ nghĩa mức cao (nhãn, vị trí, tọa độ 3D, mô tả ngữ cảnh, quyết định hành động). \\
\hline
\textbf{Bản chất chuyển đổi} & \textbf{Dữ liệu trừu tượng $\rightarrow$ Hình ảnh} & \textbf{Hình ảnh $\rightarrow$ Hình ảnh} & \textbf{Hình ảnh $\rightarrow$ Dữ liệu ngữ nghĩa} \\
\hline
\end{tabularx}
\end{table}

\subsection{Khái quát các bài toán cốt lõi trong Thị giác máy tính}

Các hệ thống thị giác máy tính hiện đại được cấu thành từ các bài toán cơ sở, trong đó mỗi bài toán được xác định rõ ràng qua mục tiêu (Goal), không gian dữ liệu đầu vào (Input) và dạng thức dữ liệu đầu ra (Output):

\textbf{1. Classification (Phân loại ảnh):}
\begin{itemize}
    \item \textbf{Mục tiêu (Goal):} Xác định và gán nhãn ngữ nghĩa ở mức toàn cảnh (Image-level) cho bức ảnh từ một tập danh mục lớp định trước.
    \item \textbf{Đầu vào (Input):} Hình ảnh tĩnh $I \in \mathbb{R}^{H \times W \times C}$ (với chiều cao $H$, chiều rộng $W$ và $C$ kênh màu).
    \item \textbf{Đầu ra (Output):} Nhãn danh mục $c \in \mathcal{C}$ hoặc vector phân phối xác suất dự đoán $P(Y = c \mid I)$ trên không gian $|\mathcal{C}|$ lớp.
\end{itemize}

\textbf{2. Detection (Phát hiện đối tượng):}
\begin{itemize}
    \item \textbf{Mục tiêu (Goal):} Đồng thời xác định vị trí không gian (Localization) và phân loại ngữ nghĩa (Classification) cho tất cả các thực thể quan tâm trong ảnh.
    \item \textbf{Đầu vào (Input):} Hình ảnh tĩnh hoặc từng khung hình video $I \in \mathbb{R}^{H \times W \times C}$.
    \item \textbf{Đầu ra (Output):} Tập hợp các hộp bao đối tượng $\mathcal{B} = \{(x_i, y_i, w_i, h_i, c_i, s_i)\}_{i=1}^N$, trong đó $(x_i, y_i, w_i, h_i)$ là tọa độ tâm và kích thước khung chữ nhật bao bọc (Bounding Box), $c_i \in \mathcal{C}$ là nhãn lớp và $s_i \in [0, 1]$ là điểm số tin cậy (Confidence Score).
\end{itemize}

\textbf{3. Recognition (Nhận dạng):}
\begin{itemize}
    \item \textbf{Mục tiêu (Goal):} Định danh cụ thể một cá thể (Instance-level identification) trong một tập danh tính/cơ sở dữ liệu đã biết (Closed-set) sau khi vùng đối tượng đã được cô lập.
    \item \textbf{Đầu vào (Input):} Vùng ảnh chứa đối tượng đã cô lập $I_{\text{crop}} \in \mathbb{R}^{h \times w \times C}$ và cơ sở dữ liệu mẫu đã đăng ký $\mathcal{D}_{\text{gallery}}$.
    \item \textbf{Đầu ra (Output):} Mã định danh cá thể cụ thể (Identity ID) $y \in \mathcal{Y}_{\text{known}}$ hoặc nhãn không xác định (\textit{Unknown}) kèm điểm số so khớp tương đồng.
\end{itemize}

\textbf{4. Tracking (Theo dõi đối tượng):}
\begin{itemize}
    \item \textbf{Mục tiêu (Goal):} Duy trì mã định danh duy nhất (Track ID) và ước lượng quỹ đạo chuyển động không gian theo thời gian của từng đối tượng qua chuỗi khung hình liên tiếp.
    \item \textbf{Đầu vào (Input):} Chuỗi khung hình video liên tục $\{I_t\}_{t=1}^T$ với $I_t \in \mathbb{R}^{H \times W \times C}$.
    \item \textbf{Đầu ra (Output):} Tập các vết quỹ đạo không gian - thời gian $\mathcal{T} = \{(\text{ID}_i, t, x_{i,t}, y_{i,t}, w_{i,t}, h_{i,t})\}$, đảm bảo tính nhất quán định danh của từng cá thể qua trục thời gian $t$.
\end{itemize}

\textbf{5. Retrieval (Truy vấn hình ảnh):}
\begin{itemize}
    \item \textbf{Mục tiêu (Goal):} Tìm kiếm và xếp hạng các hình ảnh có độ tương đồng ngữ nghĩa cao nhất với ảnh truy vấn từ một kho cơ sở dữ liệu lớn ở môi trường mở (Open-set).
    \item \textbf{Đầu vào (Input):} Ảnh truy vấn (Query Image) $I_q$ và kho cơ sở dữ liệu ảnh mẫu $\mathcal{D} = \{I_1, I_2, \dots, I_M\}$.
    \item \textbf{Đầu ra (Output):} Danh sách xếp hạng thứ tự (\textit{Ranked List}) các ảnh $\{I_{\pi(1)}, I_{\pi(2)}, \dots, I_{\pi(K)}\}$ có khoảng cách biểu diễn nhỏ nhất trong không gian nhúng đặc trưng (\textit{Embedding Space}).
\end{itemize}

\textbf{6. Counting (Đếm đối tượng):}
\begin{itemize}
    \item \textbf{Mục tiêu (Goal):} Ước lượng tổng số lượng thực thể xuất hiện trong khung cảnh, đặc biệt hiệu quả trong các tình huống có mật độ dày đặc hoặc che khuất nghiêm trọng.
    \item \textbf{Đầu vào (Input):} Hình ảnh tĩnh $I \in \mathbb{R}^{H \times W \times C}$.
    \item \textbf{Đầu ra (Output):} Số lượng thực thể ước tính $N \in \mathbb{N}$ hoặc bản đồ mật độ liên tục (Density Map) $D \in \mathbb{R}^{H \times W}$, trong đó tổng tích phân $\sum_{x,y} D(x, y)$ xấp xỉ tổng số đối tượng.
\end{itemize}

\textbf{7. Generating (Sinh dữ liệu thị giác):}
\begin{itemize}
    \item \textbf{Mục tiêu (Goal):} Khởi tạo hình ảnh hoặc video mới hoàn chỉnh, mô phỏng đúng phân phối dữ liệu thị giác thực tế $P(X)$ hoặc sinh có điều kiện theo yêu cầu.
    \item \textbf{Đầu vào (Input):} Vector nhiễu ngẫu nhiên $z \sim \mathcal{N}(0, I)$ trong không gian tiềm ẩn (Latent Space), kèm theo điều kiện dẫn đường tùy chọn (như mô tả văn bản $y_{\text{text}}$ hoặc ảnh điều kiện $I_{\text{cond}}$).
    \item \textbf{Đầu ra (Output):} Hình ảnh hoặc chuỗi video tổng hợp mới $\hat{I} \in \mathbb{R}^{H \times W \times C}$ có độ chân thực cao và thỏa mãn các ràng buộc ngữ nghĩa đầu vào.
\end{itemize}

\subsection{So sánh bài toán Nhận dạng (Recognition) và Truy vấn ảnh (Image Retrieval)}

Mặc dù đều phục vụ mục đích xác định sự tương đồng của đối tượng thị giác, Nhận dạng (Recognition) và Truy vấn ảnh (Image Retrieval) khác biệt căn bản về không gian bài toán, cơ chế tối ưu hóa hàm mất mát và khả năng mở rộng hệ thống.

\begin{table}[ht]
\centering
\small
\caption{Đối chiếu tiêu chí giữa bài toán Recognition và Image Retrieval}
\label{tab:comp_rec_retrieval}
\renewcommand{\arraystretch}{1.25}
\begin{tabularx}{\textwidth}{|p{2.3cm}|X|X|}
\hline
\textbf{Tiêu chí} & \textbf{Nhận dạng (Recognition)} & \textbf{Truy vấn hình ảnh (Image Retrieval)} \\
\hline
\textbf{Không gian bài toán} & Tập đóng (\textit{Closed-set}): số lượng lớp và danh tính được cố định trước trong quá trình huấn luyện. & Tập mở (\textit{Open-set}): tìm kiếm các mẫu và thực thể chưa từng xuất hiện trong tập huấn luyện. \\
\hline
\textbf{Mục tiêu học} & Học các siêu phẳng phân tách ranh giới giữa các lớp xác định ($P(Y=c \mid X)$). & Học không gian nhúng (\textit{Metric Learning}) sao cho các vector cùng ngữ nghĩa nằm gần nhau và khác ngữ nghĩa nằm xa nhau. \\
\hline
\textbf{Hàm mất mát} & Cross-Entropy Loss, Focal Loss. & Contrastive Loss, Triplet Loss, CosFace, ArcFace. \\
\hline
\textbf{Cơ chế suy luận} & Dự đoán phân phối xác suất lớp bằng Softmax / Argmax. & Tính toán khoảng cách (Cosine, Euclidean) kết hợp thuật toán tìm kiếm lân cận gần nhất (k-NN, ANN). \\
\hline
\textbf{Khả năng mở rộng} & Kém linh hoạt: phải huấn luyện lại hoặc tinh chỉnh (Fine-tuning) mạng khi bổ sung danh tính mới. & Rất cao: chỉ cần trích xuất vector đặc trưng và nạp vào cơ sở dữ liệu mà không cần huấn luyện lại mô hình. \\
\hline
\end{tabularx}
\end{table}

\subsection{So sánh bài toán Truy vấn ảnh (Retrieval) và Theo dõi đối tượng (Tracking)}

Sự khác biệt cốt lõi giữa Truy vấn ảnh (Retrieval) và Theo dõi đối tượng (Tracking) bắt nguồn từ yếu tố phụ thuộc thời gian, miền không gian tìm kiếm và việc kết hợp các mô hình động học:

\begin{table}[ht]
\centering
\small
\caption{Đối chiếu tiêu chí giữa bài toán Image Retrieval và Object Tracking}
\label{tab:comp_retrieval_tracking}
\renewcommand{\arraystretch}{1.25}
\begin{tabularx}{\textwidth}{|p{2.3cm}|X|X|}
\hline
\textbf{Tiêu chí} & \textbf{Truy vấn hình ảnh (Retrieval)} & \textbf{Theo dõi đối tượng (Tracking)} \\
\hline
\textbf{Bản chất dữ liệu} & Kho ảnh tĩnh phi thời gian, các mẫu dữ liệu độc lập và không có liên kết thời gian vật lý. & Chuỗi khung hình video liên tục theo thời gian thực ($t, t+1, \dots$), có tương quan chặt chẽ về không - thời gian. \\
\hline
\textbf{Phạm vi tìm kiếm} & Toàn cục (\textit{Global Search}) trên quy mô toàn bộ kho dữ liệu cơ sở $\mathcal{D}$. & Cục bộ (\textit{Local Search}) trong vùng không gian lân cận dự đoán từ khung hình liền kề. \\
\hline
\textbf{Mô hình tích hợp} & Sử dụng thuần túy mô hình đặc trưng ngoại hình tĩnh (\textit{Static Appearance Model}). & Kết hợp mô hình ngoại hình (\textit{Appearance}) và mô hình chuyển động động học (\textit{Motion Model}: Kalman Filter, Optical Flow). \\
\hline
\textbf{Dạng thức đầu ra} & Danh sách các hình ảnh xếp hạng theo độ tương đồng (\textit{Ranked List}). & Mã định danh duy nhất (\textit{Track ID}) kèm quỹ đạo tọa độ không gian liên tục theo trục thời gian. \\
\hline
\textbf{Thách thức kỹ thuật} & Tối ưu hóa chỉ mục vector quy mô lớn (\textit{Vector Indexing}) và giảm độ trễ truy vấn dữ liệu lớn. & Hiện tượng che khuất tạm thời (\textit{Occlusion}), biến dạng chuyển động nhanh (\textit{Motion Blur}), và trôi dạt định danh (\textit{ID Switch}). \\
\hline
\end{tabularx}
\end{table}

\subsection{So sánh bài toán Đếm đối tượng (Counting) và Phát hiện đối tượng (Object Detection)}

Trong các khung cảnh thông thường với mật độ thưa, bài toán đếm có thể giải quyết gián tiếp thông qua phát hiện từng đối tượng. Tuy nhiên, khi mật độ tăng cao và xuất hiện che khuất nghiêm trọng, hai phương pháp này thể hiện sự khác biệt rõ rệt:

\begin{table}[ht]
\centering
\small
\caption{Đối chiếu tiêu chí giữa bài toán Object Detection và Crowd Counting}
\label{tab:comp_detection_counting}
\renewcommand{\arraystretch}{1.25}
\begin{tabularx}{\textwidth}{|p{2.3cm}|X|X|}
\hline
\textbf{Tiêu chí} & \textbf{Phát hiện đối tượng (Detection)} & \textbf{Đếm đối tượng (Counting)} \\
\hline
\textbf{Phương pháp tiếp cận} & Định vị và hồi quy hộp bao từng cá thể độc lập (\textit{Bounding Box Regression}). & Hồi quy trường mật độ liên tục (\textit{Density Map Estimation}) trên toàn bộ khung cảnh. \\
\hline
\textbf{Dữ liệu gán nhãn} & Yêu cầu gán nhãn hộp bao chi tiết $(x, y, w, h)$ cho từng thực thể. & Chỉ yêu cầu nhãn điểm tâm (\textit{Dot/Point Annotation}) tại vị trí của từng cá thể. \\
\hline
\textbf{Độ bền vững mật độ cao} & Suy giảm nghiêm trọng: các hộp bao chồng lấn triệt tiêu lẫn nhau qua thuật toán lọc non-max (\textit{NMS}). & Bền vững cao: tổng tích phân mật độ không phụ thuộc vào việc tách biệt ranh giới pixel hay hộp bao từng cá thể. \\
\hline
\textbf{Độ phức tạp tính toán} & Tỉ lệ thuận với số lượng vùng đề xuất (Proposals), Anchors và số lượng đối tượng trong ảnh. & Cố định theo độ phân giải của bản đồ mật độ, không bị bùng nổ tính toán khi số lượng cá thể tăng cao. \\
\hline
\textbf{Dạng thức đầu ra} & Tập hợp các hộp bao rời rạc $\mathcal{B} \rightarrow N = |\mathcal{B}|$. & Bản đồ mật độ $D \in \mathbb{R}^{H \times W} \rightarrow \hat{N} = \sum_{x,y} D(x, y)$. \\
\hline
\end{tabularx}
\end{table}

\subsection{Mô hình Tạo sinh (Generative) và tương quan với các tác vụ Phân tích (Discriminative)}

Các bài toán trong học máy và thị giác máy tính được phân chia thành hai trường phái nền tảng dựa trên mục tiêu mô hình hóa phân phối xác suất và hướng biến đổi thông tin:

\begin{table}[ht]
\centering
\small
\caption{Đối chiếu tiêu chí giữa mô hình Phân tích (Discriminative) và mô hình Tạo sinh (Generative)}
\label{tab:comp_disc_gen}
\renewcommand{\arraystretch}{1.25}
\begin{tabularx}{\textwidth}{|p{2.3cm}|X|X|}
\hline
\textbf{Tiêu chí} & \textbf{Tác vụ Phân tích (Discriminative)} & \textbf{Tác vụ Tạo sinh (Generative)} \\
\hline
\textbf{Bản chất xác suất} & Học phân phối xác suất có điều kiện $P(Y \mid X)$: dự đoán nhãn hoặc quyết định $Y$ dựa trên dữ liệu $X$. & Học phân phối dữ liệu thực $P(X)$ hoặc phân phối đồng thời có điều kiện $P(X \mid Y)$. \\
\hline
\textbf{Chiều biến đổi thông tin} & Ánh xạ giảm chiều (\textit{Dimensionality Reduction}): nén dữ liệu thị giác cao chiều sang không gian nhãn tóm lược ($f: \mathcal{X} \rightarrow \mathcal{Y}$). & Ánh xạ tăng chiều/tái tạo: sinh dữ liệu thị giác chi tiết từ không gian tiềm ẩn thấp chiều ($g: \mathcal{Z} \times \mathcal{Y} \rightarrow \mathcal{X}$). \\
\hline
\textbf{Hàm mất mát chủ đạo} & Cross-Entropy Loss, Focal Loss, Smooth L1 Loss, GIoU Loss. & Adversarial Loss (Minimax), KL-Divergence, Evidence Lower Bound (ELBO), Denoising MSE. \\
\hline
\textbf{Kiến trúc tiêu biểu} & ResNet, Vision Transformer (ViT), YOLO, Faster R-CNN, Mask R-CNN. & GAN (StyleGAN), VAE, Diffusion Models (DDPM, Stable Diffusion), Autoregressive Models (DALL-E). \\
\hline
\textbf{Ứng dụng chính} & Phân loại, phát hiện khuyết tật, nhận dạng khuôn mặt, phân đoạn ảnh, hỗ trợ ra quyết định tự động. & Tổng hợp hình ảnh/video từ văn bản, tăng cường dữ liệu huấn luyện (\textit{Data Augmentation}), phục chế và nâng cao chất lượng ảnh (\textit{Super-Resolution, Inpainting}). \\
\hline
\end{tabularx}
\end{table}

\newpage
\section{Tuần 2: Trí tuệ Nhân tạo, Học máy \& Học sâu trong Khoa học Dữ liệu Thị giác}
\label{sec:week2}

\subsection{Bản chất của Trí tuệ Tự nhiên, Trí tuệ Nhân tạo và Khái niệm Thông minh}

Trong khoa học nhận thức và khoa học máy tính, việc phân định rõ các dạng thức trí tuệ là tiền đề quan trọng để xây dựng và phát triển các hệ thống tính toán:

\textbf{Trí tuệ Tự nhiên (Natural Intelligence - NI):} Là năng lực sinh học bẩm sinh của con người và các loài sinh vật bậc cao, được tôi luyện qua hàng triệu năm tiến hóa. Trí tuệ tự nhiên tích hợp toàn diện các phương thức: thu nhận cảm giác đa giác quan, tự ý thức (Self-awareness), tư duy trừu tượng, sáng tạo nghệ thuật, suy luận quan hệ nhân quả và khả năng thích ứng linh hoạt trước những biến động bất định của môi trường sống, đi kèm với đời sống cảm xúc và các chuẩn mực đạo đức xã hội.

\textbf{Trí tuệ Nhân tạo (Artificial Intelligence - AI):} Là một phân ngành của khoa học máy tính nghiên cứu các nguyên lý toán học, cấu trúc dữ liệu và giải thuật nhằm kiến tạo các hệ thống tính toán (phần mềm và phần cứng) có khả năng mô phỏng các khía cạnh hành vi thông minh của con người. Các hành vi này bao gồm năng lực học hỏi từ dữ liệu (Learning), suy luận logic (Reasoning), nhận thức thị giác (Visual Perception), xử lý ngôn ngữ tự nhiên (NLP) và lập kế hoạch ra quyết định (Decision Making).

Bảng~\ref{tab:comp_ni_ai} đối chiếu các khía cạnh nền tảng giữa Trí tuệ Tự nhiên và Trí tuệ Nhân tạo.

\begin{table}[ht]
\centering
\small
\caption{Đối chiếu nền tảng giữa Trí tuệ Tự nhiên (NI) và Trí tuệ Nhân tạo (AI)}
\label{tab:comp_ni_ai}
\renewcommand{\arraystretch}{1.3}
\begin{tabularx}{\textwidth}{|p{3.0cm}|X|X|}
\hline
\textbf{Tiêu chí đối chiếu} & \textbf{Trí tuệ Tự nhiên (NI)} & \textbf{Trí tuệ Nhân tạo (AI)} \\
\hline
\textbf{1. Cơ chất vật lý} & Mạng nơ-ron sinh học hữu cơ ($\sim 86$ tỷ nơ-ron, $10^{14}$ khớp thần kinh synapse) hoạt động thông qua tín hiệu điện sinh học và hóa sinh. & Phần cứng bán dẫn vô cơ (Silicon, GPU/TPU) xử lý thông tin thông qua các xung điện nhị phân (bit $0$ và $1$). \\
\hline
\textbf{2. Nguồn gốc \& Mục đích} & Hình thành qua quá trình chọn lọc tự nhiên và tiến hóa sinh học; mục đích cốt lõi là sinh tồn, thích nghi sinh thái và duy trì nòi giống. & Do con người thiết kế và lập trình; mục đích là công cụ tự động hóa, giải quyết bài toán kỹ thuật và phục vụ lợi ích kinh tế - xã hội. \\
\hline
\textbf{3. Cơ chế học tập} & Học hiệu quả từ số ít mẫu (Few-shot), tiếp thu qua trải nghiệm tương tác với thế giới thực và duy trì quá trình học liên tục suốt đời. & Dựa trên tối ưu hóa thống kê trên các tập dữ liệu số khổng lồ; dễ gặp hiện tượng quên tri thức cũ (Catastrophic Forgetting) khi đổi miền dữ liệu. \\
\hline
\textbf{4. Bản chất tư duy} & Suy luận quan hệ nhân quả (Causal Reasoning), trực giác trừu tượng hóa và hiểu rõ bản chất nguyên nhân (\textbf{``Tại sao''}). & So khớp mẫu và tính toán tương quan xác suất ($P(Y \mid X)$), xấp xỉ hàm toán học trả lời câu hỏi dự đoán (\textbf{``Cái gì''}). \\
\hline
\textbf{5. Ý thức \& Cảm xúc} & Có ý thức tự thân, cảm xúc thực và gắn liền với sự thấu cảm cùng trách nhiệm đạo đức. & Hoạt động dựa trên xử lý tín hiệu và tính toán số học thuần túy; phản hồi cảm xúc chỉ là kết quả mô phỏng bề mặt qua văn bản hoặc hình ảnh. \\
\hline
\textbf{6. Hiệu suất năng lượng} & Tối ưu hóa cao: toàn bộ não bộ con người chỉ tiêu thụ mức công suất sinh học ước tính khoảng $20\text{ W}$. & Tiêu hao năng lượng lớn: các trung tâm dữ liệu huấn luyện mô hình nền tảng tiêu thụ từ hàng kilowatt đến hàng megawatt điện. \\
\hline
\textbf{7. Phạm vi năng lực} & Trí tuệ tổng quát (General Intelligence): con người có thể học đồng thời nhiều kỹ năng khác nhau và chuyển giao tri thức linh hoạt. & Trí tuệ chuyên biệt (Narrow AI): thực thi vượt trội trong tác vụ hẹp được huấn luyện (cờ vây, nhận diện khuôn mặt) nhưng khó tổng quát hóa sang miền mới. \\
\hline
\end{tabularx}
\end{table}

Xét về mặt vận hành, mọi hệ thống thông minh đều thực thi một vòng lặp nhận thức khép kín gồm 3 giai đoạn liên hoàn:
\begin{enumerate}
    \item \textbf{Giai đoạn 1: Nhận biết (Perception):} Tiếp nhận các kích thích và tín hiệu vật lý từ môi trường bên ngoài. Ở con người, đó là hệ thống thụ cảm đa giác quan (thị giác, thính giác, xúc giác) với cơ chế chú ý tự nhiên chọn lọc thông tin. Ở AI, đó là các cảm biến quang học, LiDAR hoặc micro thu nhận sóng vật lý và chuyển đổi thành mảng ma trận điểm ảnh số thô $I \in \mathbb{R}^{H \times W \times C}$.
    \item \textbf{Giai đoạn 2: Hiểu (Understanding / Recognition):} Chuyển hóa dữ liệu thô thành ngữ nghĩa, cấu trúc tri thức và mô hình hóa thế giới. Con người hiểu sâu sắc về chức năng, bối cảnh văn hóa và quan hệ nhân quả của thực thể. Trong khi đó, hệ thống AI thực hiện trích xuất đặc trưng phân cấp, ánh xạ xác suất lớp $P(Y \mid X)$, định vị hộp bao và xây dựng đồ thị quan hệ cảnh (Scene Graph).
    \item \textbf{Giai đoạn 3: Thích nghi và Hành động (Action / Adaptation):} Ra quyết định, lập kế hoạch và thực thi hành động để tác động trở lại môi trường hoặc đạt tới mục tiêu tối ưu. Con người hành động dựa trên ý chí, cảm xúc và khả năng ứng biến linh hoạt trước tình huống mới. AI hành động thông qua cơ cấu chấp hành (bộ điều khiển xe tự lái, tín hiệu điều khiển relay) hoặc đề xuất giải pháp, đồng thời thích nghi bằng cách cập nhật trọng số mô hình qua giải thuật học máy.
\end{enumerate}

Từ tiến trình nhận thức trên, bản chất của khái niệm \textbf{Thông minh (Intelligence)} được xác lập dựa trên hai trụ cột cốt lõi:
\begin{itemize}
    \item \textit{Năng lực nhận biết và mô hình hóa vấn đề nhanh chóng:} Khả năng trích xuất thông tin cấu trúc và nắm bắt quy luật tiềm ẩn từ môi trường dữ liệu phức tạp, không đầy đủ hoặc có độ nhiễu cao.
    \item \textit{Năng lực thích ứng linh hoạt (Adaptability):} Khả năng tự điều chỉnh hành vi, cập nhật trọng số tri thức và đưa ra quyết định tối ưu khi môi trường xung quanh hoặc phân phối dữ liệu đầu vào biến đổi.
\end{itemize}

\subsection{Tác động của Trí tuệ Nhân tạo đối với sự phát triển của Trí tuệ Tự nhiên}

Việc ứng dụng Trí tuệ Nhân tạo vào đời sống xã hội và nghiên cứu khoa học mang lại tác động biện chứng hai chiều đối với sự phát triển của trí tuệ con người:

\textbf{1. Khía cạnh tích cực: Khuếch đại năng lực nhận thức (Cognitive Augmentation):}
\begin{itemize}
    \item \textit{Giải phóng con người khỏi các tác vụ tính toán lặp lại:} AI đảm nhận các công việc tốn thời gian nhưng mang tính quy chuẩn cao (như thống kê dữ liệu, rà soát văn bản, phân loại ảnh hàng loạt), tạo điều kiện cho con người tập trung vào tư duy sáng tạo, phản biện và hoạch định chiến lược.
    \item \textit{Mở rộng ranh giới khám phá khoa học:} AI hỗ trợ giải mã các cấu trúc quy luật phức tạp mà trực giác sinh học khó bao quát, điển hình như việc dự đoán cấu trúc protein 3D và các tương tác phân tử sinh học với AlphaFold~\citep{week2_alphafold2021,week2_alphafold2024}, tìm kiếm thiên thể trong thiên văn học hay chẩn đoán sớm tổn thương tế bào y sinh.
\end{itemize}

\textbf{2. Khía cạnh tiêu cực: Nguy cơ phụ thuộc và suy giảm nhận thức (Cognitive Offloading):}
\begin{itemize}
    \item \textit{Suy giảm phản xạ tư duy độc lập:} Sự phụ thuộc quá mức vào các đề xuất tự động của thuật toán (Automation Bias) có thể làm suy yếu khả năng tính toán độc lập, lập luận logic nền tảng và làm giảm sút năng lực ghi nhớ làm việc (Working Memory).
    \item \textit{Hiện tượng định kiến trong bong bóng thông tin (Echo Chambers):} Các thuật toán tối ưu hóa tương tác có thể cô lập người dùng trong các luồng quan điểm một chiều, làm suy giảm tính đa dạng và sự thấu cảm trong nhận thức xã hội.
\end{itemize}

Như vậy, việc ứng dụng AI không tất yếu dẫn đến sự phát triển vượt bậc hay làm thoái hóa trí tuệ tự nhiên. Hệ quả thực tế phụ thuộc vào phương thức con người định vị công nghệ: nếu xem AI là cộng sự khuếch đại nhận thức (Co-pilot) để nâng tầm tư duy thì trí tuệ tự nhiên sẽ tiến bộ; ngược lại nếu sử dụng công nghệ như một giải pháp thay thế thụ động hoàn toàn thì năng lực tư duy tự thân sẽ suy giảm.

\subsection{Phân định bản chất giữa Thông minh (Intelligence) và Thông tuệ (Wisdom)}

Trong khoa học nhận thức và triết học, ``Thông minh'' và ``Thông tuệ'' đại diện cho hai cấp độ nhận thức hoàn toàn khác biệt (Bảng~\ref{tab:comp_intelligence_wisdom}).

\begin{table}[ht]
\centering
\small
\caption{Đối chiếu toàn diện giữa khái niệm Thông minh (Intelligence) và Thông tuệ (Wisdom)}
\label{tab:comp_intelligence_wisdom}
\renewcommand{\arraystretch}{1.3}
\begin{tabularx}{\textwidth}{|p{3.0cm}|X|X|}
\hline
\textbf{Tiêu chí} & \textbf{Thông minh (Intelligence)} & \textbf{Thông tuệ / Trí huệ (Wisdom)} \\
\hline
\textbf{Định nghĩa cốt lõi} & Năng lực giải quyết bài toán và tối ưu hóa mục tiêu: thực thi công việc một cách hiệu quả và chính xác nhất (\textit{Doing things right}). & Năng lực định hướng giá trị và mục đích: lựa chọn đúng mục tiêu có ý nghĩa nhân văn để thực hiện (\textit{Doing the right things}). \\
\hline
\textbf{Đối tượng xử lý} & Dữ liệu số, cấu trúc thông tin, quy tắc logic và các không gian trạng thái toán học. & Ý nghĩa xã hội, chuẩn mực đạo đức, hạnh phúc con người và tính bền vững lâu dài. \\
\hline
\textbf{Động lực vận hành} & Tối ưu hóa hàm mục tiêu (cực tiểu hóa tổn thất hoặc cực đại hóa phần thưởng) dựa trên dữ liệu thống kê quá khứ. & Cân bằng hài hòa giữa lợi ích cá nhân và cộng đồng, gắn liền với lòng trắc ẩn (Compassion) và sự thấu cảm. \\
\hline
\textbf{Khả năng hiện thực} & Máy tính và AI có thể đạt trình độ vượt trội (Super-human) trong các tác vụ tính toán chuyên biệt. & Là đặc tính độc bản của nhận thức sinh học và trải nghiệm sống; AI hiện tại chưa thể có trí tuệ thông tuệ tự thân. \\
\hline
\end{tabularx}
\end{table}

\subsection{Giải thích Luận điểm: ``Thông minh là Khả năng Thích ứng với sự Thay đổi''}

Nhà vật lý lý thuyết Stephen Hawking từng đưa ra định nghĩa nổi tiếng về bản chất của trí tuệ:
\begin{quote}
\textit{``Intelligence is the ability to adapt to change.''} (\textit{Thông minh là khả năng thích ứng với sự thay đổi}.)
\end{quote}

Luận điểm này phản ánh nguyên lý cốt lõi của tính thông minh trong cả tiến hóa sinh học lẫn khoa học máy tính:

\textbf{1. Góc độ Trí tuệ Sinh học và Tâm lý học Nhận thức:}
Môi trường tự nhiên luôn biến động không ngừng. Một sinh vật chỉ duy trì các phản xạ cố định sẽ nhanh chóng bị đào thải khi hoàn cảnh sống thay đổi. Nhà tâm lý học Jean Piaget~\citep{week2_piaget1947} chỉ ra rằng sự phát triển nhận thức thực chất là quá trình thích ứng liên tục thông qua hai cơ chế bổ trợ: \textbf{Đồng hóa (Assimilation)}, tức quá trình tiếp nhận thông tin mới vào cấu trúc tri thức sẵn có; và \textbf{Điều ứng (Accommodation)}, tức việc tái cấu trúc mô hình tư duy để tương thích với thực tế biến đổi. Khả năng điều chỉnh linh hoạt giữa hai cơ chế này giúp sinh vật thích nghi và tồn tại.

\textbf{2. Góc độ Học máy và Thị giác Máy tính:}
Trong kỹ thuật học máy, một mô hình chỉ ghi nhớ máy móc các mẫu trong tập huấn luyện (hiện tượng Overfitting) không thể coi là thông minh. Tính thông minh của hệ thống AI được định lượng bằng năng lực thích ứng với các biến số mới:
\begin{itemize}
    \item \textit{Khả năng khái quát hóa và xử lý dữ liệu lệch phân phối (Out-of-Distribution - OOD):} Mô hình phải nhận diện chính xác đối tượng khi điều kiện chiếu sáng, góc quan sát hoặc môi trường bị biến đổi so với tập dữ liệu huấn luyện ban đầu.
    \item \textit{Thích ứng miền (Domain Adaptation) và Học chuyển giao (Transfer Learning):} Tận dụng tri thức đã học từ miền nguồn (như dữ liệu ảnh mô phỏng đồ họa 3D) để thích ứng nhanh với miền đích thực tế (ảnh camera ngoài đời thực) với lượng dữ liệu gán nhãn tối thiểu.
    \item \textit{Học tăng cường thích ứng (Reinforcement Learning):} Tác tử tự động điều chỉnh chính sách hành động $\pi(a \mid s)$ dựa trên tín hiệu phản hồi phần thưởng từ môi trường động.
    \item \textit{Học trong ngữ cảnh (In-Context Learning):} Khả năng của các Mô hình Nền tảng trong việc thực thi tác vụ mới chỉ qua một vài câu gợi ý hướng dẫn (Prompt) hoặc vài mẫu thị giác mà không cần huấn luyện lại toàn bộ mạng.
\end{itemize}

\subsection{Bản chất của Học máy: Phân tích tư tưởng Arthur Samuel (1959)}

Năm 1959, Arthur Samuel đã đưa ra định nghĩa nền tảng về Học máy:
\begin{quote}
\textit{``Machine Learning: Field of study that gives computers the ability to learn without being explicitly programmed.''}
\end{quote}

Luận điểm trên phân định ranh giới giữa hai hệ hình tính toán:
\begin{itemize}
    \item \textbf{Lập trình truyền thống (Rule-based Systems):} Lập trình viên phải hiểu rõ toàn bộ logic nghiệp vụ, sau đó viết thủ công từng dòng lệnh điều kiện (\texttt{if-else}) và các công thức giải tích. Mô hình vận hành theo sơ đồ:
    \[
    \text{Dữ liệu đầu vào (Data)} + \text{Quy tắc lập trình (Rules)} \longrightarrow \text{Kết quả đầu ra (Output)}
    \]
    Cách tiếp cận này gặp rào cản lớn trước các bài toán thị giác phi cấu trúc, bởi con người không thể lập trình tường minh hàng triệu điều kiện cho mọi biến thiên cường độ pixel của một bức ảnh.
    \item \textbf{Hệ hình Học máy (Learning-based Systems):} Thay vì lập trình quy tắc, kỹ sư thiết kế một kiến trúc mô hình tham số hóa $f_\theta(x)$ và định nghĩa một hàm mục tiêu (Loss Function) $\mathcal{L}$. Thuật toán sẽ tự động tối ưu hóa tập tham số $\theta$ dựa trên dữ liệu mẫu:
    \[
    \text{Dữ liệu đầu vào (Data)} + \text{Dữ liệu nhãn/Mục tiêu (Output)} \longrightarrow \text{Mô hình đã học ($\theta^*$)}
    \]
\end{itemize}

\subsection{Định nghĩa Hình thức về Học máy theo Tom Mitchell (1997)}

Năm 1997, Tom Mitchell đã chuẩn hóa lý thuyết học máy bằng định nghĩa toán học chặt chẽ:
\begin{quote}
\textit{``A computer program is said to learn from experience $\mathcal{E}$ with respect to some class of tasks $\mathcal{T}$ and performance measure $\mathcal{P}$, if its performance at tasks in $\mathcal{T}$, as measured by $\mathcal{P}$, improves with experience $\mathcal{E}$.''}
\end{quote}

Định nghĩa này cấu thành từ bộ ba $(\mathcal{T}, \mathcal{E}, \mathcal{P})$:
\begin{itemize}
    \item \textbf{Nhiệm vụ $\mathcal{T}$ (Task):} Bài toán mục tiêu mà hệ thống cần thực hiện. Ví dụ trong thị giác máy tính: Phát hiện vị trí các khuyết tật trên bề mặt vi mạch $I \in \mathbb{R}^{H \times W \times C}$.
    \item \textbf{Trải nghiệm $\mathcal{E}$ (Experience):} Tập hợp dữ liệu lịch sử hoặc các tương tác mà mô hình được tiếp cận. Ví dụ: Tập dữ liệu $\mathcal{D} = \{(x_i, y_i)\}_{i=1}^N$ gồm $100,000$ hình ảnh vi mạch được gán nhãn vị trí lỗi chuẩn bởi chuyên gia kiểm thử.
    \item \textbf{Độ đo hiệu năng $\mathcal{P}$ (Performance Measure):} Thước đo định lượng toán học dùng để đánh giá chất lượng thực thi tác vụ. Ví dụ: Chỉ số độ chính xác trung bình $\text{mAP}@0.5$ trên tập kiểm thử độc lập $\mathcal{D}_{\text{test}}$.
\end{itemize}

Một thuật toán được xác nhận là đang học khi chất lượng thực thi tăng đơn điệu theo lượng dữ liệu trải nghiệm:
\begin{equation}
\frac{\partial \mathcal{P}(\mathcal{T})}{\partial \mathcal{E}} > 0 \quad \text{hoặc} \quad \lim_{|\mathcal{E}| \to \infty} \mathcal{L}(\theta; \mathcal{E}) \to \min
\end{equation}

\subsection{Tiến trình Chuyển tiếp từ Học máy sang Học sâu: Giới hạn và Động lực Đột phá}

Sự phát triển của thị giác máy tính gắn liền với bước chuyển dịch từ các thuật toán học máy cổ điển sang các kiến trúc mạng nơ-ron sâu:

\textbf{1. Quá trình chuyển dịch công nghệ:}
Trong thời kỳ học máy truyền thống (trước năm 2012), quy trình xử lý tách rời làm hai giai đoạn độc lập: (1) Trích xuất đặc trưng thủ công (Hand-crafted Feature Extraction) do con người thiết kế như SIFT~\citep{lowe2004distinctive}, HOG~\citep{dalal2005histograms} hay LBP; (2) Đưa các vector đặc trưng này vào bộ phân loại nông như SVM~\citep{cortes1995support} hoặc Random Forest~\citep{breiman2001random}. Ngược lại, kỷ nguyên học sâu tích hợp toàn bộ quy trình thành một mạng duy nhất theo cơ chế học biểu diễn đầu-cuối (End-to-End Representation Learning), ánh xạ trực tiếp từ pixel thô sang dự đoán ngữ nghĩa.

\textbf{2. Những hạn chế căn bản của Học máy truyền thống:}
\begin{itemize}
    \item \textit{Nút thắt cổ chai từ đặc trưng thủ công:} Các bộ trích xuất thủ công chỉ nắm bắt được các mẫu hình học cục bộ đơn giản (cạnh, góc, kết cấu xám), dễ mất tính bất biến trước các biến dạng phối cảnh lớn, che khuất hoặc thay đổi tư thế phức tạp.
    \item \textit{Lời nguyền số chiều (Curse of Dimensionality):} Khi xử lý ảnh có độ phân giải cao, không gian đặc trưng bùng nổ khiến các thuật toán dựa trên ranh giới phẳng hoặc khoảng cách hình học bị suy giảm hiệu năng nghiêm trọng.
    \item \textit{Sự bão hòa dung lượng mô hình:} Khi tăng quy mô dữ liệu từ hàng vạn lên hàng triệu mẫu, hiệu năng của các mô hình nông nhanh chóng đạt trần do dung lượng biểu diễn (Model Capacity) bị giới hạn.
\end{itemize}

\textbf{3. Động lực thúc đẩy sự ra đời của Học sâu:}
\begin{itemize}
    \item \textit{Tự động học biểu diễn phân cấp (Hierarchical Feature Abstraction):} Mạng nơ-ron sâu tự động trừu tượng hóa thông tin theo từng tầng: từ đường biên cục bộ ở tầng đầu, đến các bộ phận hình thể ở tầng giữa, và ngữ nghĩa đối tượng hoàn chỉnh ở tầng sâu nhất.
    \item \textit{Định luật mở rộng quy mô (Scaling Laws):} Hiệu năng của học sâu liên tục tăng trưởng tương ứng với kích thước dữ liệu và số lượng tham số khi được hỗ trợ bởi năng lực tính toán GPU/TPU song song quy mô lớn.
    \item \textit{Năng lực biểu diễn đa phương thức:} Khả năng liên kết biểu diễn đồng thời ảnh, video, văn bản và âm thanh vào cùng một không gian nhúng chung.
\end{itemize}

\subsection{Nền tảng của Học sâu: Vì sao Nguyên lý Đơn giản tạo nên Khả năng Biểu diễn Vượt trội?}

Về mặt toán học, mạng nơ-ron sâu được xây dựng từ các phép toán nền tảng:
\begin{enumerate}
    \item Phép biến đổi affine tuyến tính: $\mathbf{z} = \mathbf{W}\mathbf{x} + \mathbf{b}$.
    \item Phép kích hoạt phi tuyến cục bộ từng phần tử: $\mathbf{a} = \sigma(\mathbf{z})$ (với $\sigma$ là hàm ReLU, GELU, Sigmoid).
    \item Quá trình tối ưu hóa bằng giải thuật lan truyền ngược (Backpropagation) dựa trên quy tắc đạo hàm chuỗi (Chain Rule) và phương pháp hạ dốc độ dốc ngẫu nhiên (SGD~\citep{robbins1951stochastic}).
\end{enumerate}

Nguyên lý toán học này mang lại khả năng biểu diễn vượt trội nhờ ba cơ chế nền tảng:
\begin{itemize}
    \item \textbf{Định lý Xấp xỉ Phổ quát (Universal Approximation Theorem):} Một mạng nơ-ron chứa các tầng ẩn phi tuyến có khả năng xấp xỉ bất kỳ hàm số liên tục nào với độ chính xác tùy ý. Khi xếp chồng nhiều tầng sâu, cấu trúc phân cấp cho phép mô hình biểu diễn các hàm logic phức tạp với số lượng tham số tiết kiệm hơn theo cấp số nhân so với mạng nông.
    \item \textbf{Cấu trúc thiên lệch cảm ứng (Inductive Bias):} Kiến trúc tích chập (Convolution) tận dụng tính chất bất biến tịnh tiến và tính tương quan cục bộ của hình ảnh tự nhiên, giúp thu hẹp không gian tìm kiếm tham số một cách tối ưu. Tương tự, cơ chế tự chú ý (Self-Attention) trong Vision Transformer mô hình hóa quan hệ phụ thuộc toàn cục linh hoạt.
    \item \textbf{Hình học tối ưu trong không gian tham số cao chiều:} Trong không gian hàng triệu đến hàng tỷ chiều, các điểm dừng tối ưu chủ yếu là các điểm yên ngựa (Saddle Points) mà giải thuật tối ưu hóa có thể vượt qua, giúp mô hình hội tụ về các vùng nghiệm có khả năng khái quát hóa cao.
\end{itemize}

\subsection{Mô hình 5 Cấp độ Nhận thức và Dữ liệu: Ứng dụng Giám sát An ninh Trường học}

Quá trình chuyển hóa dữ liệu từ dạng tín hiệu vật lý thô sang hành động có giá trị nhân văn được mô hình hóa qua \textbf{5 Cấp độ Nhận thức}:
\begin{enumerate}
    \item \textbf{Level 1: Unfiltered Data (Dữ liệu thô chưa lọc):} Tín hiệu điểm ảnh, sóng âm, dữ liệu cảm biến chưa qua xử lý.
    \item \textbf{Level 2: Pattern Information (Thông tin mẫu):} Các mẫu hình học, đặc trưng thị giác, vector nhúng được trích xuất có cấu trúc.
    \item \textbf{Level 3: Predictability Knowledge (Tri thức có tính dự đoán):} Nhận định danh tính, xu hướng, mối liên hệ nhân quả và quy luật lịch sử.
    \item \textbf{Level 4: Choice Intelligence (Trí thông minh ra quyết định):} Năng lực lựa chọn phương án hành động tối ưu trong các kịch bản thực thi.
    \item \textbf{Level 5: Compassion Wisdom (Thông tuệ thấu cảm):} Sự đánh giá giá trị nhân văn, thấu hiểu hoàn cảnh và đưa ra giải pháp nhân bản.
\end{enumerate}

Bảng~\ref{tab:matrix_5_levels_school} phân tích ma trận ứng dụng 5 cấp độ này vào bài toán \textbf{Hệ thống Giám sát An ninh \& Điểm danh Khuôn mặt tại Trường học}, chỉ rõ từng Task, Cấp độ nhận thức tương ứng, Dữ liệu đầu vào, Dữ liệu đầu ra và Nội dung chi tiết trong từng ô.

\begin{table}[ht]
\centering
\scriptsize
\setlength{\tabcolsep}{3.5pt}
\caption{Ma trận phân tích 5 cấp độ nhận thức ứng dụng trong Hệ thống Giám sát An ninh Trường học}
\label{tab:matrix_5_levels_school}
\renewcommand{\arraystretch}{1.3}
\begin{tabularx}{\textwidth}{|p{2.1cm}|p{2.0cm}|p{2.5cm}|p{2.4cm}|X|}
\hline
\textbf{Task hệ thống} & \textbf{Cấp độ nhận thức} & \textbf{Dữ liệu đầu vào (Input)} & \textbf{Dữ liệu đầu ra (Output)} & \textbf{Nội dung chi tiết \& Ví dụ minh họa trong ô} \\
\hline
\textbf{Task 1: Thu nhận luồng hình ảnh} & \textbf{Level 1: Unfiltered Data} & Luồng tín hiệu video RTSP từ camera giám sát tại cổng trường & Mảng ma trận byte điểm ảnh thô $I_t \in \mathbb{R}^{H \times W \times 3}$ & \textbf{[Mức 1 - Dữ liệu thô]:} Thu nhận 30 khung hình/giây chứa mọi thông tin nền (bóng râm, cây cối, xe cộ, ánh sáng đổi) mà máy tính chưa phân biệt được có học sinh hay không. \\
\hline
\textbf{Task 2: Phát hiện \& Trích xuất khuôn mặt} & \textbf{Level 2: Pattern Information} & Khung ảnh thô $I_t$ & Hộp bao khuôn mặt $\mathcal{B}$ và Vector nhúng đặc trưng $\mathbf{e} \in \mathbb{R}^{512}$ & \textbf{[Mức 2 - Thông tin mẫu]:} Thuật toán RetinaFace định vị khuôn mặt và trích xuất vector 512 chiều biểu diễn cấu trúc hình học khoảng cách mắt, mũi, miệng. \\
\hline
\textbf{Task 3: Định danh \& Đối chiếu lịch sử} & \textbf{Level 3: Predictability Knowledge} & Vector nhúng $\mathbf{e}$ kèm cơ sở dữ liệu học sinh $\mathcal{D}_{\text{student}}$ & Danh tính học sinh kèm độ tin cậy, lịch sử chuyên cần & \textbf{[Mức 3 - Tri thức dự đoán]:} Xác định khuôn mặt là ``Em An - Lớp 10A1'' với độ khớp 99.4\%; nhận diện học sinh đi trễ 3 ngày liên tiếp, dự báo nguy cơ vi phạm nền nếp. \\
\hline
\textbf{Task 4: Điều khiển cổng \& Cảnh báo tự động} & \textbf{Level 4: Choice Intelligence} & Kết quả định danh và quy tắc an ninh trường học & Tín hiệu relay mở cổng / Gửi thông báo phụ huynh / Chuông báo & \textbf{[Mức 4 - Trí thông minh lựa chọn]:} Nếu là học sinh: tự động mở chốt cửa và gửi thông báo điểm danh cho phụ huynh; nếu phát hiện người lạ chưa đăng ký: tự động khóa cổng và báo bảo vệ. \\
\hline
\textbf{Task 5: Đánh giá hoàn cảnh \& Hỗ trợ nhân văn} & \textbf{Level 5: Compassion Wisdom} & Nhật ký đi trễ, biểu cảm khuôn mặt mệt mỏi và dữ liệu gia cảnh & Đề xuất tư vấn tâm lý, hỗ trợ hoàn cảnh khó khăn & \textbf{[Mức 5 - Thông tuệ thấu cảm]:} Không trừ điểm thi đua máy móc; hệ thống ghi nhận học sinh đi trễ kèm biểu hiện kiệt sức, đề xuất giáo viên chủ nhiệm thăm hỏi gia cảnh và hỗ trợ tâm lý học đường. \\
\hline
\end{tabularx}
\end{table}

\subsection{Mô hình Giải pháp 3 Tầng trong Khoa học Dữ liệu Thị giác}

Để xây dựng một hệ thống thị giác máy tính hoàn chỉnh phục vụ thực tiễn, kiến trúc kỹ thuật được chuẩn hóa thành \textbf{Mô hình Giải pháp 3 Tầng} (Bảng~\ref{tab:three_tier_cv_architecture}):
\begin{enumerate}
    \item \textbf{Tầng 1: Kỹ thuật Cốt lõi / Nền tảng Thuật toán (Foundational Technologies):} Bao gồm các mô hình toán học, kiến trúc mạng nơ-ron nền tảng (CNN, Vision Transformer), hàm mất mát, thuật toán tối ưu (AdamW~\citep{loshchilov2019decoupled}, SGD~\citep{robbins1951stochastic}) và giải thuật hình học thị giác (Epipolar, Homography).
    \item \textbf{Tầng 2: Tác vụ Thị giác Cơ sở (Core Vision Tasks):} Các khối chức năng thị giác cơ bản giải quyết các bài toán độc lập như Phân loại ảnh, Phát hiện hộp bao, Nhận dạng khuôn mặt, Ước lượng chiều sâu, Tách mặt nạ phân đoạn hay Theo dõi quỹ đạo.
    \item \textbf{Tầng 3: Hệ thống Ứng dụng Nghiệp vụ Hoàn chỉnh (End-to-End Business Applications):} Tích hợp các khối chức năng ở Tầng 2 cùng các ràng buộc nghiệp vụ, giao diện người dùng, cơ sở dữ liệu và phần cứng điều khiển để tạo thành sản phẩm hoàn chỉnh giải quyết bài toán thực tế.
\end{enumerate}

\begin{table}[ht]
\centering
\small
\caption{Cấu trúc Mô hình Giải pháp 3 Tầng trong Thị giác máy tính}
\label{tab:three_tier_cv_architecture}
\renewcommand{\arraystretch}{1.3}
\begin{tabularx}{\textwidth}{|p{2.8cm}|X|X|}
\hline
\textbf{Phân tầng} & \textbf{Thành phần cốt lõi} & \textbf{Ví dụ minh họa cụ thể} \\
\hline
\textbf{Tầng 1: Kỹ thuật cốt lõi} & Kiến trúc mạng (ResNet~\citep{he2016deep}, Swin-ViT~\citep{liu2021swin}, YOLO-backbone~\citep{redmon2016you}), Hàm mất mát (ArcFace, InfoNCE, GIoU~\citep{rezatofighi2019generalized}), Thư viện tối ưu (PyTorch, TensorRT). & Mạng ResNet-50~\citep{he2016deep} kết hợp hàm mất mát ArcFace tối ưu hóa bằng AdamW~\citep{loshchilov2019decoupled} trên tập dữ liệu chuẩn. \\
\hline
\textbf{Tầng 2: Tác vụ thị giác} & Các bài toán cơ sở: Face Detection, Facial Landmark Alignment, Re-Identification, Crowd Density Estimation. & Phát hiện 68 điểm mốc khuôn mặt và trích xuất vector đặc trưng khuôn mặt 512 chiều. \\
\hline
\textbf{Tầng 3: Ứng dụng nghiệp vụ} & Hệ thống hoàn chỉnh tích hợp cơ sở dữ liệu, phân quyền, cảnh báo thời gian thực và tự động hóa quy trình nghiệp vụ. & Hệ thống Kiosk tự động Check-in hành khách tại sân bay tích hợp cổng kiểm soát an ninh xuất nhập cảnh. \\
\hline
\end{tabularx}
\end{table}

\subsection{Khái quát các Nhóm Ứng dụng Cơ bản của Thị giác Máy tính}

Thị giác máy tính hiện đại được phân bổ thành 5 nhóm ứng dụng công nghiệp trọng yếu (Bảng~\ref{tab:cv_application_domains}).

\begin{table}[ht]
\centering
\small
\caption{Tổng hợp 5 nhóm lĩnh vực ứng dụng cơ bản của Thị giác máy tính}
\label{tab:cv_application_domains}
\renewcommand{\arraystretch}{1.3}
\begin{tabularx}{\textwidth}{|p{3.2cm}|X|X|}
\hline
\textbf{Lĩnh vực ứng dụng} & \textbf{Các tác vụ Tầng 2 tích hợp} & \textbf{Hệ thống sản phẩm Tầng 3 tiêu biểu} \\
\hline
\textbf{1. Giám sát an ninh \& Đô thị thông minh} & Detection, Tracking, License Plate Recognition, Crowd Counting. & Hệ thống camera phạt nguội giao thông tự động; Giám sát trật tự đô thị và an ninh sân bay. \\
\hline
\textbf{2. Y tế số \& Chẩn đoán hình ảnh} & Semantic Segmentation, 3D Voxel Reconstruction, Anomaly Detection. & Phần mềm tự động khoanh vùng khối u trên ảnh chụp CT/MRI; Phát hiện sớm polyp nội soi tiêu hóa. \\
\hline
\textbf{3. Tự hành \& Robot công nghiệp} & 3D Object Detection, Depth Estimation, Visual SLAM, Optical Flow. & Hệ thống xe tự lái; Cánh tay robot tự động gắp thả linh kiện và kiểm tra lỗi hàn vi mạch. \\
\hline
\textbf{4. Bán lẻ \& Thương mại điện tử} & Image Retrieval, Fine-grained Classification, Pose Estimation. & Tìm kiếm sản phẩm bằng hình ảnh trên sàn TMĐT; Cửa hàng bán lẻ tự động không người thu ngân. \\
\hline
\textbf{5. Điện ảnh, CGI \& Thực tế ảo (AR/VR)} & 3D Face Reconstruction, Motion Capture, Generative Video Inpainting. & Kỹ xảo điện ảnh kỹ thuật số (VFX); Tạo Avatar Metahuman; Bộ lọc hiệu ứng khuôn mặt AR trên smartphone. \\
\hline
\end{tabularx}
\end{table}

\subsection{Phân tích Chuyên sâu: Ứng dụng Trí tuệ Nhân tạo trong Ngành Điện ảnh}

Ứng dụng AI trong phân tích trải nghiệm khán giả điện ảnh (AI for Cinema Audience Analysis) là ví dụ điển hình về sự phối hợp giữa Tầng 1, Tầng 2 và Tầng 3 nhằm tối ưu hóa giá trị nghệ thuật và thương mại (Bảng~\ref{tab:ai_cinema_pipeline}):

\begin{table}[ht]
\centering
\small
\caption{Quy trình phân tích trải nghiệm khán giả điện ảnh ứng dụng AI}
\label{tab:ai_cinema_pipeline}
\renewcommand{\arraystretch}{1.3}
\begin{tabularx}{\textwidth}{|p{3.0cm}|X|X|}
\hline
\textbf{Giai đoạn phân tích} & \textbf{Công nghệ \& Tác vụ thị giác (Tầng 1 \& 2)} & \textbf{Giá trị ứng dụng nghiệp vụ điện ảnh (Tầng 3)} \\
\hline
\textbf{1. Nhận diện khuôn mặt \& Cảm xúc} & Camera hồng ngoại trong rạp tối $\rightarrow$ Facial Landmark Detection $\rightarrow$ Phân loại biểu cảm (Vui mừng, Sợ hãi, Bất ngờ, Thất vọng) theo hệ thống mã hóa hành vi nét mặt FACS. & Đo lường chính xác mức độ kích thích cảm xúc của người xem tại từng cảnh phim; kiểm chứng xem phân đoạn phim có đạt hiệu ứng mong muốn (ví dụ cảnh hài có làm khán giả cười hay không). \\
\hline
\textbf{2. Phân tích hành vi \& Hướng nhìn} & Head Pose Estimation, Eye Gaze Tracking $\rightarrow$ Action Recognition (che mắt, ngả người ra sau, quay sang người bên cạnh, xem điện thoại). & Đánh giá mức độ tập trung (Engagement Level) của khán giả; phát hiện các phân đoạn gây buồn ngủ hoặc nhàm chán khiến khán giả phân tâm. \\
\hline
\textbf{3. Đồng bộ Timeline \& Tối ưu kịch bản} & Tích hợp chuỗi thời gian cảm xúc khán giả đồng bộ với biểu đồ Timeline dựng phim (Cut-by-cut synchronization). & Hỗ trợ đạo diễn và nhà sản xuất trong khâu chiếu thử nghiệm (Test Screening): cắt gọt các phân đoạn lê thê, tinh chỉnh nhịp độ (Pacing), dự báo doanh thu phòng vé và định hướng chiến dịch Marketing. \\
\hline
\end{tabularx}
\end{table}

\subsection{Ý nghĩa Thực tiễn của Bài toán Đếm Đám đông ở Tầng Ứng dụng Nghiệp vụ}

Ở Tầng 2 (Tác vụ kỹ thuật), bài toán Đếm đám đông (Crowd Counting) giải quyết việc ánh xạ hình ảnh $I \in \mathbb{R}^{H \times W \times C}$ sang bản đồ mật độ liên tục $D(x, y)$ để ước lượng số thực thể $\hat{N} = \sum_{x,y} D(x, y)$.

Khi được triển khai lên \textbf{Tầng 3 (Hệ thống Ứng dụng Nghiệp vụ Thực tiễn)}, bài toán này mang lại giá trị quản trị và an toàn xã hội to lớn:
\begin{itemize}
    \item \textbf{Phòng ngừa thảm họa giẫm đạp (Stampede Prevention):} Tại các sự kiện lễ hội, nhà ga tàu điện ngầm hoặc sân vận động, hệ thống liên tục tính toán mật độ cục bộ ($\text{người}/\text{m}^2$). Khi mật độ vượt ngưỡng an toàn ($>4\text{ người}/\text{m}^2$), hệ thống tự động kích hoạt cảnh báo nguy cơ tắc nghẽn cục bộ để lực lượng an ninh kịp thời phân luồng trước khi xảy ra tai nạn.
    \item \textbf{Tối ưu hóa phân bổ nguồn lực công cộng:} Dữ liệu mật độ thời gian thực giúp cơ quan quản lý điều phối tần suất phương tiện giao thông công cộng, bố trí nhân viên cứu hộ và mở thêm cửa thoát hiểm linh hoạt.
    \item \textbf{Phân tích hành vi tiêu dùng và quy hoạch thương mại:} Đếm lưu lượng người đi bộ qua các gian hàng trong trung tâm thương mại giúp xác định các vị trí tập trung đông khách để định giá mặt bằng cho thuê và tối ưu hóa vị trí đặt biển quảng cáo.
\end{itemize}

\subsection{Phân tích Hệ thống Tái tạo Khuôn mặt 3D theo Mô hình Giải pháp 3 Tầng}

Hệ thống tái tạo khuôn mặt 3D (3D Face Reconstruction) từ ảnh 2D đơn góc nhìn là minh chứng tiêu biểu cho kiến trúc phân tầng hoàn chỉnh từ lý thuyết toán học đến ứng dụng thực tế (Bảng~\ref{tab:three_tier_3d_face}).

\begin{table}[ht]
\centering
\small
\caption{Phân tích hệ thống Tái tạo Khuôn mặt 3D theo Mô hình Giải pháp 3 Tầng}
\label{tab:three_tier_3d_face}
\renewcommand{\arraystretch}{1.3}
\begin{tabularx}{\textwidth}{|p{2.8cm}|X|X|}
\hline
\textbf{Phân tầng} & \textbf{Nhiệm vụ \& Thành phần kỹ thuật} & \textbf{Chi tiết giải pháp công nghệ} \\
\hline
\textbf{Tầng 1: Kỹ thuật cốt lõi} & Cơ sở mô hình toán hình học 3D, kiến trúc mạng nơ-ron và hàm tối ưu khả vi (Differentiable Rendering). & Mô hình tham số hình học 3D Morphable Model (3DMM: BFM, FLAME); Mạng Transformer trích xuất hệ số hình dạng $\boldsymbol{\alpha}$, biểu cảm $\boldsymbol{\beta}$ và tư thế $\mathbf{p}$; Hàm mất mát Photometric Consistency Loss, Landmark Reprojection Loss, và Perceptual Identity Loss. \\
\hline
\textbf{Tầng 2: Tác vụ thị giác cơ sở} & Các khối trích xuất thông tin trung gian từ ảnh 2D sang không gian 3D. & Phát hiện 68/468 điểm mốc 3D (Dense Facial Landmarks); Ước lượng trường độ sâu đơn sắc (Monocular Depth Estimation); Tách bản đồ phản xạ ánh sáng (Albedo/Texture UV Map) và góc chiếu sáng (Spherical Harmonics). \\
\hline
\textbf{Tầng 3: Ứng dụng nghiệp vụ hoàn chỉnh} & Tích hợp vào các phần mềm ứng dụng phục vụ người dùng cuối. & \begin{itemize}[leftmargin=*,noitemsep,topsep=0pt]
    \item \textbf{Y tế \& Thẩm mỹ:} Mô phỏng kết quả phẫu thuật chỉnh hình hàm mặt trước khi phẫu thuật thực tế.
    \item \textbf{Giải trí \& Game:} Tự động tạo Digital Avatar 3D cá nhân hóa có khả năng cử động khớp thời gian thực trong game và không gian thực tế ảo.
    \item \textbf{Bảo mật sinh trắc học:} Xác thực khuôn mặt 3D chống giả mạo (Face Anti-spoofing) ngăn chặn hành vi vượt rào bằng ảnh in 2D hoặc màn hình kỹ thuật số.
\end{itemize} \\
\hline
\end{tabularx}
\end{table}

\newpage
\section{Tuần 3: Ứng dụng Đồ họa Máy tính và Đề tài Nghiên cứu Sinh Đồ họa Vector (Text-to-Vector)}
\label{sec:week3}

\subsection{Bản chất và Mục tiêu của Đồ họa Máy tính}

Đồ họa Máy tính (Computer Graphics) là một phân ngành nền tảng của khoa học máy tính, nghiên cứu các giải pháp phần cứng và phần mềm nhằm mô phỏng, biểu diễn và tổng hợp thế giới thực hoặc các thế giới trừu tượng dưới dạng thông tin thị giác và âm thanh kỹ thuật số~\citep{hearn2004computer,foley1996computer}.

Về mặt lịch sử và lý thuyết mô hình hóa, quá trình phát triển của đồ họa máy tính phản ánh tiến trình nâng bậc không gian biểu diễn toán học từ thấp lên cao:
\begin{equation}
\text{2D } (x, y) \longrightarrow \text{3D } (x, y, z) \longrightarrow \text{4D } (x, y, z, t) \longrightarrow \text{5D } (x, y, z, t, s)
\end{equation}
Trong chuỗi tiến trình này:
\begin{itemize}
    \item \textbf{Không gian 2 chiều (2D):} Mô hình hóa các thực thể phẳng thông qua tọa độ $(x, y)$, đặt nền móng cho việc xử lý văn bản, sơ đồ kỹ thuật số, giao diện người dùng và nghệ thuật đồ họa phẳng.
    \item \textbf{Không gian 3 chiều (3D):} Bổ sung trục độ sâu $z$, cho phép định nghĩa các khối hình học lập thể, bề mặt vật thể thực trong không gian Euclide thông qua lưới đa giác (Polygon Mesh) hoặc bề mặt tham số (Parametric Surfaces).
    \item \textbf{Không gian 4 chiều (4D):} Tích hợp thêm chiều thời gian $t$, chuyển đổi từ mô hình hình học tĩnh sang hoạt họa chuyển động (Computer Animation) và mô phỏng các hiện tượng động học vật lý thời gian thực.
    \item \textbf{Không gian 5 chiều (5D):} Bổ sung chiều ngữ cảnh, thang đo kích thước hoặc trạng thái cảm biến $s$ (State / Scale / Semantics), phục vụ các hệ thống mô phỏng đa tỉ lệ, đa phương thức trong thực tế ảo tương tác (VR/AR) và các bản sao số kỹ thuật số (Digital Twins).
\end{itemize}

\subsection{Phân tích Ứng dụng của Đồ họa Máy tính trong 7 Lĩnh vực Trọng yếu}

Đồ họa máy tính đóng vai trò công nghệ cốt lõi trong toàn bộ nền kinh tế số. Bảng~\ref{tab:cg_applications_detailed} phân tích chi tiết ứng dụng của đồ họa máy tính trên 7 lĩnh vực trọng yếu được giảng dạy trong học phần.

\begin{table}[ht]
\centering
\small
\caption{Tổng hợp và đối chiếu ứng dụng của Đồ họa Máy tính trong 7 lĩnh vực công nghiệp cốt lõi}
\label{tab:cg_applications_detailed}
\renewcommand{\arraystretch}{1.3}
\begin{tabularx}{\textwidth}{|p{3.2cm}|X|X|}
\hline
\textbf{Lĩnh vực ứng dụng} & \textbf{Đặc trưng kỹ thuật đồ họa cốt lõi} & \textbf{Hệ thống sản phẩm thực tế tiêu biểu} \\
\hline
\textbf{1. Tự động hóa văn phòng (Office Automation)} & Trực quan hóa dữ liệu thống kê (Bar, Pie, Line charts); giao diện người dùng đồ họa (GUI); biểu diễn biểu đồ luồng nghiệp vụ. & Bộ ứng dụng Microsoft Office, Google Workspace; phần mềm xây dựng Dashboard thông minh như Tableau, Microsoft PowerBI. \\
\hline
\textbf{2. Xuất bản để bàn (Desktop Publishing - DTP)} & Xử lý typography vector tham số hóa (TrueType, OpenType); phân tách màu CMYK; quản lý dàn trang độ phân giải cao; chuẩn PostScript/PDF. & Adobe InDesign, QuarkXPress, Scribus; hệ thống biên soạn tài liệu khoa học chất lượng cao \LaTeX. \\
\hline
\textbf{3. Thiết kế \& Chế tạo trợ giúp máy tính (CAD-CAM)} & Mô hình hóa hình học chính xác (NURBS, CSG, B-rep); phân tích ứng suất phần tử hữu hạn (FEA); xuất mã G-code cho máy tiện/phay CNC. & AutoCAD, Dassault Systèmes SolidWorks, CATIA, PTC Creo; phần mềm thiết kế vi mạch Altium Designer. \\
\hline
\textbf{4. Nghệ thuật \& Quảng cáo (Art + Advertising)} & Chỉnh sửa ảnh raster đa lớp; thiết kế nhận diện thương hiệu vector; nghệ thuật minh họa số (Digital Painting); Typography 3D. & Adobe Photoshop, Adobe Illustrator, CorelDRAW; phần mềm vẽ kỹ thuật số Procreate, Clip Studio Paint. \\
\hline
\textbf{5. Điều khiển quy trình (Process Control / SCADA)} & Giao diện tương tác người máy (HMI) thời gian thực; sơ đồ động học hệ thống (Mimic Diagrams); hiển thị cảnh báo trực quan đa cấp độ. & Hệ thống SCADA giám sát nhà máy nhiệt điện, trạm biến áp, dây chuyền lắp ráp ô tô tự động, lọc dầu của Siemens, Schneider Electric. \\
\hline
\textbf{6. Giải trí (Entertainment: Phim \& Games)} & Kết xuất quang học (Ray Tracing, Global Illumination); mô phỏng vật lý hạt (khói, lửa, nước); tạo hình nhân vật và diễn hoạt 3D thời gian thực. & Unreal Engine 5, Unity; phần mềm 3D chuyên nghiệp Autodesk Maya, Blender; phim hoạt hình của Pixar, Walt Disney. \\
\hline
\textbf{7. Giáo dục \& Đào tạo (Education)} & Mô phỏng hình học không gian tương tác; phòng thí nghiệm ảo (Virtual Labs); buồng mô phỏng lái máy bay/tàu thủy (Flight Simulators). & Hệ thống mô phỏng bay CAE Flight Simulator; phần mềm toán học tương tác GeoGebra; phòng học ảo 3D mô phỏng giải phẫu học cơ thể người. \\
\hline
\end{tabularx}
\end{table}

\subsection{Đề tài Nghiên cứu Chuyên sâu: Sinh Đồ họa Vector từ Ngôn ngữ Tự nhiên (Text-to-Vector Graphics Generation)}

Trong khuôn khổ môn học Đồ họa Máy tính, nhóm nghiên cứu lựa chọn đề tài \textbf{``Sinh Đồ họa Vector từ Ngôn ngữ Tự nhiên'' (Text-to-Vector Graphics Generation, đặc biệt hướng tới định dạng chuẩn Scalable Vector Graphics - SVG)}. Dưới đây là phân tích toàn diện về bản chất học thuật và giá trị ứng dụng thực tiễn của đề tài.

\subsubsection{Bản chất Học thuật và Thách thức Nghiên cứu}

Trong khi lĩnh vực AI tạo sinh đã đạt được những bước tiến vượt bậc với việc sinh ảnh điểm ảnh (Raster Image Generation) nhờ các mạng khuếch tán (Diffusion Models) như Stable Diffusion hay Midjourney, việc sinh đồ họa vector lại đặt ra một thách thức hoàn toàn khác biệt về mặt bản thể luận:

\textbf{1. Sự lệch pha giữa Trí tuệ Thị giác Nhân tạo và Tư duy Thiết kế Con người:}
Các mô hình học sâu truyền thống tư duy trên không gian lưới điểm ảnh cố định (Fixed Pixel Grids), nơi bức ảnh được biểu diễn bằng một ma trận tensor có kích thước xác định $I \in \mathbb{R}^{H \times W \times C}$. Ngược lại, đồ họa vector và tư duy của các nhà thiết kế đồ họa vận hành dựa trên các thực thể hình học tham số (Geometric Primitives) như đường cong Bézier bậc ba, hình tròn, đa giác và các lớp đối tượng phân cấp ngữ nghĩa (Hierarchical Layers). Sự chuyển dịch từ miền điểm ảnh rời rạc sang miền đường cong tham số liên tục là một bài toán khó do tính chất phi khả vi (Non-differentiable) của các trình raster hóa truyền thống.

\textbf{2. Toàn cảnh 5 hướng tiếp cận nghiên cứu học thuật hiện đại:}
Để giải quyết bài toán Text-to-Vector, cộng đồng khoa học quốc tế trong giai đoạn 2020 đến nay đã phát triển 5 nhánh tiếp cận chủ đạo:
\begin{enumerate}
    \item \textbf{Tối ưu hóa lặp qua Trình kết xuất khả vi (Optimization-based Synthesis via Differentiable Rendering):} Khởi đầu từ việc Li et al. công bố DiffVG~\citep{li2020diffvg}, cho phép truyền ngược đạo hàm từ pixel về tọa độ điểm điều khiển Bézier. CLIPDraw~\citep{frans2022clipdraw} sử dụng không gian nhúng CLIP để tối ưu nét vẽ. Tiếp đó, VectorFusion~\citep{jain2023vectorfusion} áp dụng Score Distillation Sampling (SDS) từ Stable Diffusion, và SVGDreamer~\citep{xing2024svgdreamer} nâng cấp lên Variational Particle Score Distillation (VPSD) kết hợp phân rã chú ý chéo (Cross-Attention) để cải thiện độ sắc nét của hình vẽ vector.
    \item \textbf{Biểu diễn Ẩn Thần kinh và Không gian Tiềm ẩn Đường cong (Neural Implicit \& Latent Path Representations):} Nhằm khắc phục tình trạng đường cong bị răng cưa, tự cắt nhau và biến dạng khi tối ưu hóa điểm tự do, các phương pháp như T2V-NPR sử dụng VAE để học không gian tiềm ẩn của các đường cong chuẩn mực, trong khi NIVeL và NeuralSVG sử dụng mạng nơ-ron truyền thẳng (MLP) làm hàm biểu diễn ẩn, kết hợp kỹ thuật Nested Dropout để ép mạng học thứ tự lớp hình học từ lớn đến nhỏ.
    \item \textbf{Mô hình Ngôn ngữ Lớn \& Đa phương thức Sinh mã Trực tiếp (Autoregressive Code Generation with LLMs \& MLLMs):} Coi tệp SVG như một ngôn ngữ hình thức dạng XML, các nghiên cứu mới nhất như StarVector~\citep{rodriguez2025starvector} (CVPR 2025) và LLM4SVG~\citep{xing2025llm4svg} (CVPR 2025) huấn luyện các mô hình đa phương thức trên hàng triệu tệp SVG thực tế, thiết kế các token ngữ nghĩa SVG chuyên biệt để sinh trực tiếp mã SVG sạch, giữ đúng các thẻ cơ sở như \texttt{<circle>}, \texttt{<rect>}, \texttt{<polygon>} chỉ trong vài chục giây suy luận.
    \item \textbf{Lai ghép Ngôn ngữ và Thị giác (Hybrid LLM-Template + Image Diffusion Refinement):} Điển hình là Chat2SVG (CVPR 2025), giao nhiệm vụ kiến trúc sư phân rã bố cục ngữ nghĩa cho LLM (như Claude 3.5 Sonnet) để tạo khung mẫu, sau đó dùng mô hình khuếch tán ảnh và SAM để điêu khắc chi tiết vi mô cho từng thành phần.
    \item \textbf{Mô phỏng Quy trình Tư duy Thiết kế Phân lớp bằng Diffusion Transformer (Process-Centric Layered Generation via DiT):} Nghiên cứu LayerTracer (ICCV 2025) học trực tiếp quy trình vẽ từng lớp của họa sĩ thông qua kiến trúc DiT (FLUX.1) trên bố cục hình rắn (Serpentine Layout), tạo ra bản thiết kế phân lớp hoàn toàn độc lập và khớp với tư duy con người.
\end{enumerate}

\subsubsection{Ý nghĩa và Giá trị Ứng dụng Thực tiễn}

Định dạng đồ họa Scalable Vector Graphics (SVG) sở hữu những ưu thế kỹ thuật vượt trội không thể thay thế trong ngành công nghiệp đồ họa hiện đại:
\begin{itemize}
    \item \textbf{Độ phân giải độc lập (Resolution Independence):} Hình ảnh vector có thể phóng to đến kích thước vô hạn mà hoàn toàn không bị hiện tượng vỡ hạt (Pixelation), đảm bảo chất lượng hiển thị sắc nét tuyệt đối từ màn hình đồng hồ thông minh, điện thoại, máy tính màn hình Retina/4K đến các biển quảng cáo tấm lớn ngoài trời.
    \item \textbf{Dung lượng siêu nhẹ và tối ưu hóa Web:} Tập tin SVG thực chất là văn bản mã nguồn XML có cấu trúc toán học gọn nhẹ (thường chỉ vài Kilobytes), nhẹ hơn hàng chục đến hàng trăm lần so với ảnh raster (PNG, JPEG), giúp giảm tải băng thông và tăng tốc độ tải trang web.
    \item \textbf{Tính khả biên tập và tái cấu trúc (High Editability):} Khác với ảnh raster là một khối pixel hợp nhất không thể tách rời, đồ họa vector lưu trữ từng đối tượng độc lập. Người dùng có thể dễ dàng nhấp chọn từng hình khối, thay đổi màu sắc, dịch chuyển vị trí, chỉnh sửa độ cong hoặc gán chuyển động hoạt hình bằng CSS/JavaScript mà không cần vẽ lại từ đầu.
    \item \textbf{Tự động hóa chu trình sản xuất công nghiệp và thiết kế UI/UX:} Ứng dụng mô hình Text-to-Vector cho phép các nhà thiết kế và lập trình viên tự động tạo lập hàng loạt biểu tượng (Icons), hình minh họa đồ họa (Illustrations), sơ đồ giao diện người dùng và tệp cắt kỹ thuật số (Laser cutting, CNC) chỉ từ các câu lệnh mô tả ngắn, rút ngắn thời gian thiết kế từ nhiều giờ xuống còn vài giây.
\end{itemize}

\newpage
\section{Tuần 4: Hệ Tọa độ, Cơ chế Chiếu và Giải thuật Vẽ Đoạn thẳng, Đường tròn}
\label{sec:week4}

\subsection{Bản chất và Phân định Ba Hệ Tọa độ trong Đồ họa Máy tính (WCS, UCS, DCS)}

Trong quy trình kết xuất và xử lý đồ họa máy tính, việc quản lý không gian hình học được phân định thông qua ba hệ tọa độ nền tảng, phản ánh các giai đoạn chuyển hóa dữ liệu từ mô hình trừu tượng sang màn hình hiển thị vật lý:

\textbf{1. Hệ Tọa độ Thế giới Thực (World Coordinate System, WCS):}
WCS là hệ tọa độ Decartes tuyệt đối toàn cục được dùng làm chuẩn tham chiếu chung cho toàn bộ không gian cảnh đồ họa. Trong hệ tọa độ này, mọi đối tượng, nguồn sáng, góc nhìn camera và các điều kiện biên môi trường đều được xác định vị trí tương đối với nhau. Đơn vị đo trong WCS mang ý nghĩa vật lý thực tế (như mét, milimét, kilômét hoặc inch), cho phép mô hình hóa các vật thể theo đúng tỉ lệ kích thước tự nhiên.

\textbf{2. Hệ Tọa độ Người dùng / Mô hình hóa (User Coordinate System, UCS / Modeling Coordinate System, MCS):}
UCS là hệ tọa độ cục bộ độc lập do người thiết kế hoặc lập trình viên tự định nghĩa cho từng đối tượng hình học riêng lẻ. Việc sử dụng UCS giúp quá trình tạo dựng mô hình trở nên thuận tiện và trực quan, bởi gốc tọa độ thường được đặt tại tâm đối xứng hoặc một đỉnh đặc trưng của vật thể. Người dùng có thể tự do tịnh tiến, xoay hoặc thay đổi hướng của các trục tọa độ cục bộ này trong quá trình thiết kế. Để đưa một vật thể từ không gian mô hình hóa UCS vào không gian chung WCS, hệ thống đồ họa áp dụng phép biến đổi mô hình (Modeling Transformation) thông qua ma trận biến đổi affine:
\begin{equation}
\mathbf{P}_{\text{WCS}} = \mathbf{M}_{\text{model}} \cdot \mathbf{P}_{\text{UCS}}
\end{equation}

\textbf{3. Hệ Tọa độ Thiết bị / Màn hình Hiển thị (Device Coordinate System, DCS / Screen Coordinate System):}
DCS là hệ tọa độ 2D rời rạc tương ứng trực tiếp với lưới điểm ảnh vật lý của thiết bị xuất (màn hình hiển thị hoặc máy in). Trục tọa độ trong DCS được định nghĩa bằng các giá trị số nguyên:
\begin{equation}
x_{\text{DCS}} \in \{0, 1, \dots, X_{\max}\}, \quad y_{\text{DCS}} \in \{0, 1, \dots, Y_{\max}\}
\end{equation}
Gốc tọa độ $(0, 0)$ của DCS thường được quy ước đặt tại góc trên bên trái (phổ biến trong hệ thống cửa sổ đồ họa máy tính) hoặc góc dưới bên trái của màn hình. 

Để chuyển đổi một vùng quan sát trong thế giới thực sang một khung nhìn trên màn hình, đồ họa máy tính thực thi phép biến đổi Cửa sổ sang Khung nhìn (Window-to-Viewport Transformation, WinToView). Nếu cửa sổ quan sát trong thế giới thực được giới hạn bởi $[x_{w\min}, x_{w\max}] \times [y_{w\min}, y_{w\max}]$ và khung nhìn màn hình được giới hạn bởi $[x_{v\min}, x_{v\max}] \times [y_{v\min}, y_{v\max}]$, tọa độ hiển thị $(x_v, y_v)$ được ánh xạ tuyến tính theo công thức:
\begin{equation}
x_v = x_{v\min} + (x_w - x_{w\min}) \frac{x_{v\max} - x_{v\min}}{x_{w\max} - x_{w\min}}, \quad y_v = y_{v\min} + (y_w - y_{w\min}) \frac{y_{v\max} - y_{v\min}}{y_{w\max} - y_{w\min}}
\end{equation}

\subsection{Cơ sở Lý thuyết: Mô phỏng Thế giới Thực thông qua Biểu diễn Vector Kích thước Nhỏ}

Trong đồ họa máy tính, một nghịch lý thú vị thường được đặt ra: \textit{Tại sao máy tính có thể mô phỏng chính xác thế giới thực rộng lớn khi kích thước lưu trữ của các vector dữ liệu lại vô cùng nhỏ gọn?}

Bản chất của vấn đề nằm ở hai nguyên lý khoa học nền tảng:

\textbf{1. Trừu tượng hóa Toán học và Biểu diễn Tham số:}
Đồ họa vector không lưu trữ hàng tỷ hạt vật chất hay hàng triệu điểm ảnh rời rạc của thế giới thực. Thay vào đó, nó trừu tượng hóa thế giới thành các quy luật toán học hình học tham số (Parametric Geometric Models). Ví dụ, một mặt cầu trong không gian thực có kích thước bán kính hàng nghìn kilômét chỉ cần được biểu diễn bằng một bộ 4 số thực: tọa độ tâm $(x_c, y_c, z_c)$ và bán kính $r$. Tương tự, một bề mặt khí động học phức tạp của cánh máy bay được mô tả chính xác tuyệt đối bằng một lưới tập hợp vài chục điểm điều khiển của mặt cong NURBS (Non-Uniform Rational B-Splines). Dung lượng của các tham số số học này chỉ chiếm vài chục đến vài trăm byte nhớ, nhưng mang đầy đủ thông tin hình học giải tích chuẩn xác.

\textbf{2. Tính Bất biến về Thang đo và Ánh xạ Tỉ lệ Hình học (Scale Invariance):}
Mô hình toán học hoàn toàn độc lập với kích thước vật lý của màn hình hiển thị. Nhờ phép biến đổi tỉ lệ (Scaling Transformation) thông qua ma trận đại số:
\begin{equation}
\mathbf{S}(s_x, s_y, s_z) = \begin{bmatrix} s_x & 0 & 0 & 0 \\ 0 & s_y & 0 & 0 \\ 0 & 0 & s_z & 0 \\ 0 & 0 & 0 & 1 \end{bmatrix}
\end{equation}
hệ thống đồ họa có thể co giãn thế giới mô phỏng về bất kỳ thang đo kích thước nào. Cùng một công thức vector, hệ thống có thể phóng to hàng triệu lần để quan sát cấu trúc tinh thể vi mô của phân tử nano, hoặc thu nhỏ hàng tỷ lần để mô phỏng quỹ đạo chuyển động của các thiên thể trong hệ ngân hà trên màn hình máy tính có kích thước chỉ vài chục inch.

\subsection{Cơ chế Mô phỏng Đối tượng Thực 3 Chiều trên Màn hình Hiển thị 2 Chiều}

Màn hình máy tính về bản chất vật lý là một bề mặt phẳng hai chiều. Để tái tạo chân thực cảm giác không gian ba chiều có chiều sâu cho mắt người, đồ họa máy tính triển khai một chuỗi kỹ thuật kết xuất liên hoàn (Viewing Pipeline)~\citep{hearn2004computer}:

\textbf{1. Phép Chiếu Hình học (Geometric Projections):}
Ánh xạ các điểm từ không gian 3 chiều $(x, y, z)$ xuống mặt phẳng quan sát 2 chiều $(x_p, y_p)$:
\begin{itemize}
    \item \textit{Phép chiếu Song song (Parallel / Orthographic Projection):} Các tia chiếu song song với nhau. Phép chiếu này bảo toàn tính song song của các đường thẳng và giữ nguyên tỉ lệ kích thước thực của đối tượng, không phụ thuộc vào khoảng cách tới người quan sát. Phương pháp này là tiêu chuẩn bắt buộc trong các bản vẽ kỹ thuật cơ khí và kiến trúc xây dựng (CAD).
    \item \textit{Phép chiếu Phối cảnh (Perspective Projection):} Các tia chiếu cùng hội tụ về một điểm gốc gọi là tâm chiếu (Center of Projection, tương đương vị trí mắt người quan sát hoặc ống kính máy ảnh). Phép chiếu phối cảnh tạo ra hiện tượng biến dạng xa gần tự nhiên (Foreshortening): vật thể càng ở xa tâm chiếu thì kích thước hình chiếu trên mặt phẳng chiếu càng nhỏ, và các đường thẳng song song trong không gian sẽ hội tụ tại một điểm tụ (Vanishing Point). Ma trận chiếu phối cảnh chuẩn hóa với khoảng cách tiêu cự $d$ có dạng:
    \begin{equation}
    x_p = \frac{x}{z / d}, \quad y_p = \frac{y}{z / d}
    \end{equation}
\end{itemize}

\textbf{2. Thuật toán Khử Mặt Khuất (Hidden Surface Removal, HSR):}
Trong không gian thực, các vật thể phía trước sẽ che khuất các vật thể phía sau. Đồ họa máy tính sử dụng thuật toán Đệm độ sâu (Z-Buffer / Depth-Buffer Algorithm). Tại mỗi pixel $(x, y)$ trên màn hình, hệ thống duy trì giá trị độ sâu $z$ nhỏ nhất (gần người quan sát nhất). Khi một đa giác mới được raster hóa, giá trị $z_{\text{new}}$ của nó sẽ được so sánh với $z_{\text{buffer}}(x, y)$; nếu $z_{\text{new}} < z_{\text{buffer}}(x, y)$, điểm ảnh mới sẽ được cập nhật màu sắc và ghi đè giá trị độ sâu mới vào bộ đệm.

\textbf{3. Mô hình Chiếu sáng và Tô bóng (Illumination and Shading):}
Để bề mặt vật thể thể hiện được độ cong và chất liệu thực tế, hệ thống áp dụng các mô hình quang học mô phỏng tương tác giữa ánh sáng và bề mặt vật liệu, tiêu biểu là Mô hình Chiếu sáng Phong (Phong Illumination Model):
\begin{equation}
I = I_a k_a + I_d k_d (\mathbf{L} \cdot \mathbf{N}) + I_s k_s (\mathbf{R} \cdot \mathbf{V})^n
\end{equation}
trong đó $I_a k_a$ là thành phần ánh sáng môi trường (Ambient), $I_d k_d (\mathbf{L} \cdot \mathbf{N})$ là ánh sáng khuếch tán Lambertian (Diffuse), và $I_s k_s (\mathbf{R} \cdot \mathbf{V})^n$ là ánh sáng phản xạ gương tạo vệt sáng bóng (Specular Highlight). Sau đó, kỹ thuật tô bóng Gouraud Shading (nội suy màu sắc qua các đỉnh) hoặc Phong Shading (nội suy vector pháp tuyến $\mathbf{N}$ qua từng điểm ảnh) được áp dụng để loại bỏ cảm giác đa giác sắc cạnh thô ráp.

\textbf{4. Ánh xạ Vân Bề mặt (Texture Mapping):}
Dán các bức ảnh 2D (Textures) lên bề mặt các đa giác 3D kết hợp với kỹ thuật Bump Mapping hoặc Normal Mapping để mô phỏng độ nhám, nếp nhăn và kết cấu chi tiết của vật liệu (gỗ, đá, da, kim loại) mà không làm tăng số lượng đa giác hình học cần xử lý.

\subsection{Phân tích Toàn diện Giải thuật Vẽ Đoạn thẳng và Khảo sát Độ dốc}

Bài toán vẽ đoạn thẳng trên lưới điểm ảnh số yêu cầu xác định tập hợp các pixel nguyên $(x_k, y_k)$ sao cho đường gấp khúc nối các pixel này xấp xỉ tốt nhất đoạn thẳng giải tích nối hai điểm mút $A(x_a, y_a)$ và $B(x_b, y_b)$, đồng thời đảm bảo tính liên tục không gian trong lân cận 8 pixel.

Phương trình tường minh của đoạn thẳng trong mặt phẳng 2D có dạng:
\begin{equation}
y = m \cdot x + b
\end{equation}
trong đó hệ số góc (độ dốc) $m$ và hằng số dịch chuyển $b$ được xác định bởi:
\begin{equation}
m = \frac{\Delta y}{\Delta x} = \frac{y_b - y_a}{x_b - x_a}, \quad b = y_a - m \cdot x_a
\end{equation}

Từ định nghĩa vi sai $\Delta y = m \cdot \Delta x$, ta có mối quan hệ độ lớn biến thiên:
\begin{equation}
|\Delta y| = |m| \cdot |\Delta x|
\end{equation}

Biểu thức này dẫn đến hai trường hợp cơ bản quyết định giải thuật tạo điểm ảnh:

\textbf{Trường hợp 1: Độ dốc tuyệt đối nhỏ hơn hoặc bằng 1 ($|m| \le 1 \iff |\Delta y| \le |\Delta x|$):}
Trong trường hợp này, chiều dài biến thiên theo trục hoành $x$ lớn hơn hoặc bằng chiều dài biến thiên theo trục tung $y$. Nếu ta chọn bước nhảy đơn vị theo trục hoành là $\Delta x = \pm 1$, khi đó bước nhảy tương ứng của tung độ thỏa mãn:
\begin{equation}
|\Delta y| = |m| \cdot 1 \le 1
\end{equation}
Điều này có nghĩa là mỗi khi $x$ tăng đúng 1 pixel, giá trị thực của $y$ biến thiên một khoảng không vượt quá 1 pixel. Nhờ đó, tại mỗi bước, pixel kế tiếp chỉ dịch chuyển sang cột bên cạnh và tung độ chỉ thay đổi tối đa 1 đơn vị, đảm bảo đoạn thẳng vẽ ra liên tục, không bao giờ xuất hiện hiện tượng gián đoạn hay thủng lỗ theo phương ngang. Do đó, $x$ được chọn làm biến độc lập tăng đều, và tung độ tại bước kế tiếp được tính theo công thức:
\begin{equation}
x_{k+1} = x_k + 1, \quad y_{k+1} = y_k + m
\end{equation}

\textbf{Trường hợp 2: Độ dốc tuyệt đối lớn hơn 1 ($|m| > 1 \iff |\Delta y| > |\Delta x|$):}
Nếu trong trường hợp này ta vẫn tiếp tục chọn $\Delta x = \pm 1$, bước nhảy tung độ sẽ là $|\Delta y| = |m| > 1$. Khi đó, giá trị $y$ sẽ nhảy cóc qua nhiều pixel, tạo thành các khoảng trống gián đoạn nghiêm trọng trên màn hình theo phương thẳng đứng. Để khắc phục triệt để, ta phải hoán đổi vai trò của hai trục: chọn trục tung $y$ làm biến độc lập tăng đều với $\Delta y = \pm 1$. Khi đó, bước nhảy tương ứng theo trục hoành thỏa mãn:
\begin{equation}
|\Delta x| = \frac{1}{|m|} \cdot |\Delta y| = \frac{1}{|m|} < 1
\end{equation}
Lúc này, mỗi khi $y$ tăng đúng 1 pixel, $x$ chỉ biến thiên một khoảng nhỏ hơn 1 pixel, đảm bảo tính liên tục tuyệt đối trong lân cận 8 pixel theo phương thẳng đứng:
\begin{equation}
y_{k+1} = y_k + 1, \quad x_{k+1} = x_k + \frac{1}{m}
\end{equation}

Phương pháp này chính là nền tảng của giải thuật DDA (Digital Differential Analyzer). Tuy nhiên, DDA đòi hỏi các phép tính số thực dấu phẩy động và hàm làm tròn \texttt{round()} tại mọi bước lặp, gây suy giảm tốc độ kết xuất đồ họa.

\subsection{Cơ sở Giới hạn Bài toán \texorpdfstring{$0 < m < 1, x_a < x_b$}{0 < m < 1, xa < xb} và Cơ chế Suy diễn Toàn diện}

Trong nghiên cứu thuật toán vẽ đoạn thẳng, các nhà khoa học luôn giới hạn bài toán khảo sát ban đầu về trường hợp chuẩn tắc:
\begin{equation}
0 < m < 1 \quad \text{và} \quad x_a < x_b
\end{equation}
tương ứng với góc nghiêng $\theta \in (0^\circ, 45^\circ)$ trong Bát phân thứ nhất (Octant 1).

\textbf{Lý do Giới hạn Bài toán:}
Việc giới hạn vào khoảng này giúp tối giản hóa công thức toán học và cấu trúc rẽ nhánh của giải thuật. Khi $x_a < x_b$, quá trình quét luôn tiến từ trái sang phải ($\Delta x = x_b - x_a > 0$). Khi $0 < m < 1$, tại mỗi bước hoành độ luôn tăng đúng 1 đơn vị ($x_{k+1} = x_k + 1$), và hệ thống chỉ cần đưa ra quyết định nhị phân duy nhất giữa hai pixel lân cận: pixel nằm ngang $(x_k + 1, y_k)$ hoặc pixel nằm chéo lên $(x_k + 1, y_k + 1)$.

\textbf{Cơ chế Suy diễn cho Toàn bộ các Trường hợp còn lại:}
Mọi đoạn thẳng bất kỳ trong không gian 2D đều có thể quy về trường hợp chuẩn tắc thông qua các phép biến đổi đối xứng hình học và xử lý điều kiện biên:

\begin{enumerate}
    \item \textbf{Đảm bảo chiều quét từ trái sang phải ($x_a < x_b$):} Nếu $x_a > x_b$, ta chỉ việc hoán đổi hai điểm mút: gán $A \leftrightarrow B$. Khi đó đoạn thẳng được vẽ từ $B$ sang $A$ với cùng tập hợp các pixel nhưng đảm bảo $\Delta x > 0$.
    \item \textbf{Các trường hợp suy biến đặc biệt:}
    \begin{itemize}
        \item Đoạn thẳng nằm ngang ($m = 0 \iff y_a = y_b$): Giữ nguyên $y = y_a$, tăng dần $x$ từ $x_a$ đến $x_b$ và vẽ pixel $(x, y_a)$.
        \item Đoạn thẳng đứng ($m \to \infty \iff x_a = x_b$): Giữ nguyên $x = x_a$, tăng dần $y$ từ $\min(y_a, y_b)$ đến $\max(y_a, y_b)$ và vẽ pixel $(x_a, y)$.
        \item Đoạn thẳng nghiêng đúng $45^\circ$ ($|m| = 1$): Cả $x$ và $y$ cùng tăng hoặc giảm đồng thời 1 đơn vị ở mỗi bước lặp ($x_{k+1} = x_k + 1, y_{k+1} = y_k \pm 1$).
    \end{itemize}
    \item \textbf{Suy diễn cho 8 Bát phân tổng quát (Octants):}
    Dựa vào dấu của $\Delta x, \Delta y$ và so sánh độ lớn giữa $|\Delta x|$ và $|\Delta y|$:
    \begin{itemize}
        \item Nếu $|\Delta y| \le |\Delta x|$: Trục $x$ tăng đều từng bước $1$ ($s_x = \text{sign}(\Delta x) = 1$). Nếu $m > 0$, $y$ tăng về phía dương ($s_y = 1$); nếu $m < 0$, $y$ biến thiên về phía âm ($s_y = -1$).
        \item Nếu $|\Delta y| > |\Delta x|$: Hoán đổi vai trò hai trục tọa độ ($x \leftrightarrow y$). Trục $y$ tăng đều từng bước $1$ ($s_y = \text{sign}(\Delta y) = 1$), và trục $x$ được cập nhật dựa trên hàm quyết định tương ứng.
    \end{itemize}
\end{enumerate}

\subsection{Cơ chế Khử Phép tính Số thực của Thuật toán Bresenham Đoạn thẳng}

Năm 1965, Jack E. Bresenham đã đề xuất giải thuật kinh điển mang tên ông~\citep{bresenham1965algorithm}, loại bỏ hoàn toàn các phép toán số thực và hàm làm tròn của thuật toán DDA, chuyển toàn bộ quá trình tính toán về các phép cộng và trừ số nguyên thuần túy.

Xét đoạn thẳng trong bát phân thứ nhất với $0 < m < 1$ và $x_a < x_b$. Giả sử tại bước thứ $k$, thuật toán đã xác định pixel cần vẽ là $(x_k, y_k)$. Ở bước kế tiếp $k+1$, hoành độ chắc chắn là $x_{k+1} = x_k + 1$. Ta cần lựa chọn giữa hai pixel ứng viên:
\begin{equation}
\text{Pixel dưới: } (x_k + 1, y_k) \quad \text{hoặc} \quad \text{Pixel trên: } (x_k + 1, y_k + 1)
\end{equation}

Tung độ thực tế của điểm nằm trên đoạn thẳng tại hoành độ $x_k + 1$ là:
\begin{equation}
y = m(x_k + 1) + b
\end{equation}

Khoảng cách thẳng đứng từ tung độ thực tế $y$ đến pixel dưới và pixel trên lần lượt là:
\begin{align}
d_1 &= y - y_k = m(x_k + 1) + b - y_k \\
d_2 &= (y_k + 1) - y = y_k + 1 - [m(x_k + 1) + b]
\end{align}

Hiệu số giữa hai khoảng cách này phản ánh pixel nào nằm gần đoạn thẳng thực tế hơn:
\begin{align}
d_1 - d_2 &= [m(x_k + 1) + b - y_k] - [y_k + 1 - m(x_k + 1) - b] \nonumber \\
&= 2m(x_k + 1) - 2y_k + 2b - 1
\end{align}
Thay $m = \frac{\Delta y}{\Delta x}$ vào biểu thức:
\begin{equation}
d_1 - d_2 = 2 \frac{\Delta y}{\Delta x}(x_k + 1) - 2y_k + 2b - 1
\end{equation}

Đại lượng này vẫn chứa phép chia số thực $\frac{\Delta y}{\Delta x}$. Nhận thấy rằng $\Delta x > 0$, Bresenham đã nhân toàn bộ biểu thức trên với $\Delta x$ để định nghĩa hàm tham số quyết định (Decision Parameter) $p_k$:
\begin{equation}
p_k = \Delta x (d_1 - d_2) = 2\Delta y \cdot x_k - 2\Delta x \cdot y_k + C
\end{equation}
trong đó hằng số $C = 2\Delta y + \Delta x(2b - 1)$.

Quy tắc chọn pixel dựa trên dấu của $p_k$:
\begin{itemize}
    \item Nếu $p_k < 0 \iff d_1 < d_2$: Đoạn thẳng nằm gần pixel dưới hơn, do đó ta chọn pixel $(x_k + 1, y_k)$.
    \item Nếu $p_k \ge 0 \iff d_1 \ge d_2$: Đoạn thẳng nằm gần pixel trên hơn (hoặc cách đều), ta chọn pixel $(x_k + 1, y_k + 1)$.
\end{itemize}

Để thiết lập công thức tính truy hồi, ta xét giá trị $p_{k+1}$ tại bước kế tiếp:
\begin{equation}
p_{k+1} = 2\Delta y \cdot x_{k+1} - 2\Delta x \cdot y_{k+1} + C
\end{equation}
Lấy hiệu số giữa $p_{k+1}$ và $p_k$:
\begin{align}
p_{k+1} - p_k &= 2\Delta y (x_{k+1} - x_k) - 2\Delta x (y_{k+1} - y_k) \nonumber \\
&= 2\Delta y (1) - 2\Delta x (y_{k+1} - y_k)
\end{align}
Do đó, công thức truy hồi cập nhật tham số quyết định có dạng:
\begin{equation}
p_{k+1} = p_k + 2\Delta y - 2\Delta x (y_{k+1} - y_k)
\end{equation}

Từ đây, ta có hai trường hợp cập nhật cụ thể:
\begin{itemize}
    \item Nếu $p_k < 0$: Chọn $y_{k+1} = y_k \implies y_{k+1} - y_k = 0$, ta cập nhật:
    \begin{equation}
    p_{k+1} = p_k + 2\Delta y
    \end{equation}
    \item Nếu $p_k \ge 0$: Chọn $y_{k+1} = y_k + 1 \implies y_{k+1} - y_k = 1$, ta cập nhật:
    \begin{equation}
    p_{k+1} = p_k + 2\Delta y - 2\Delta x
    \end{equation}
\end{itemize}

Vì các điểm mút ban đầu có tọa độ nguyên, nên $\Delta x, \Delta y, 2\Delta y$ và $2\Delta y - 2\Delta x$ đều là các hằng số nguyên. Toàn bộ thuật toán hoàn toàn không sử dụng bất kỳ phép tính số thực hay phép chia nào, mà chỉ gồm phép cộng, trừ số nguyên và phép nhân 2 (có thể tối ưu bằng phép dịch bit sang trái \texttt{<< 1}), mang lại hiệu năng tính toán cực cao trên phần cứng.

\subsection{Chứng minh Công thức Khởi tạo Giá trị \texorpdfstring{$p_0$}{p0}}

Để hoàn thiện giải thuật truy hồi Bresenham, ta cần tính giá trị ban đầu $p_0$ tại bước $k = 0$:

Điểm mút đầu tiên $(x_0, y_0) = (x_a, y_a)$ nằm chính xác trên đường thẳng giải tích, do đó thỏa mãn:
\begin{equation}
y_0 = m \cdot x_0 + b \implies b = y_0 - m \cdot x_0 = y_0 - \frac{\Delta y}{\Delta x} x_0
\end{equation}

Tại bước đầu tiên, ta cần quyết định pixel cho cột $x_1 = x_0 + 1$. Áp dụng công thức tổng quát của tham số quyết định:
\begin{equation}
p_0 = \Delta x [2m(x_0 + 1) - 2y_0 + 2b - 1]
\end{equation}
Thay $m = \frac{\Delta y}{\Delta x}$ và $b = y_0 - \frac{\Delta y}{\Delta x} x_0$ vào biểu thức:
\begin{align}
p_0 &= \Delta x \left[ 2 \frac{\Delta y}{\Delta x}(x_0 + 1) - 2y_0 + 2\left(y_0 - \frac{\Delta y}{\Delta x} x_0\right) - 1 \right] \nonumber \\
&= \Delta x \left[ 2 \frac{\Delta y}{\Delta x} x_0 + 2 \frac{\Delta y}{\Delta x} - 2y_0 + 2y_0 - 2 \frac{\Delta y}{\Delta x} x_0 - 1 \right] \nonumber \\
&= \Delta x \left[ 2 \frac{\Delta y}{\Delta x} - 1 \right]
\end{align}
Nhân phân phối $\Delta x$ vào bên trong dấu ngoặc vuông:
\begin{equation}
p_0 = 2\Delta y - \Delta x
\end{equation}

Đây là công thức khởi tạo chuẩn xác của thuật toán vẽ đoạn thẳng Bresenham, chứng minh rằng $p_0$ luôn luôn là một giá trị nguyên.

\subsection{Cơ sở Vẽ Hình tròn trên Lưới Tọa độ Nguyên và Tính Đối xứng 8 Phần}

Phương trình đại số của đường tròn tâm tại gốc tọa độ $O(0, 0)$ với bán kính $r$ có dạng:
\begin{equation}
x^2 + y^2 = r^2
\end{equation}
Hàm đường tròn được định nghĩa tổng quát bằng:
\begin{equation}
f_{\text{circle}}(x, y) = x^2 + y^2 - r^2
\end{equation}
Vị trí tương đối của một điểm tọa độ bất kỳ $(x, y)$ so với đường biên đường tròn được phân loại chặt chẽ qua dấu của hàm:
\begin{equation}
f_{\text{circle}}(x, y) = \begin{cases}
< 0, & \text{nếu } (x, y) \text{ nằm phía trong đường tròn} \\
= 0, & \text{nếu } (x, y) \text{ nằm chính xác trên đường biên đường tròn} \\
> 0, & \text{nếu } (x, y) \text{ nằm phía ngoài đường tròn}
\end{cases}
\end{equation}

\textbf{Nguyên lý Đối xứng 8 Phần (8-Way Symmetry):}
Đường tròn tâm $O(0, 0)$ sở hữu tính đối xứng hoàn hảo qua trục hoành $y = 0$, trục tung $x = 0$ và hai đường phân giác $y = x$, $y = -x$. Do đó, toàn bộ chu vi đường tròn được chia thành 8 cung bát phân bằng nhau. 

Nếu thuật toán tính toán được tập hợp các pixel trên cung bát phân thứ hai (từ điểm đỉnh $(0, r)$ ở góc $90^\circ$ xuống giao điểm với đường phân giác tại góc $45^\circ$, nơi $x = y$), ta có thể tự động suy ra 7 điểm còn lại trên các cung khác chỉ bằng phép đổi dấu và hoán vị tọa độ:
\begin{equation}
(x, y), \quad (y, x), \quad (-x, y), \quad (-y, x), \quad (-x, -y), \quad (-y, -x), \quad (x, -y), \quad (y, -x)
\end{equation}
Sau đó, nếu đường tròn có tâm thực tế tại $C(x_c, y_c)$, ta chỉ cần thực hiện phép tịnh tiến bằng cách cộng thêm $(x_c, y_c)$ vào từng tọa độ đã tính.

\subsection{Phân tích Tính Liên tục 8-Lân cận: Vì sao \texorpdfstring{$x$}{x} phải Tăng đều trong Cung \texorpdfstring{$[90^\circ, 45^\circ]$}{[90, 45] độ}}

Trong cung bát phân khảo sát (từ điểm đỉnh $(0, r)$ đến điểm giao $x = y$), câu hỏi đặt ra là: \textit{Để đảm bảo đường cong liên tục không bị thủng lỗ trong lân cận 8 pixel, ta phải chọn biến $x$ hay $y$ tăng đều?}

Lấy đạo hàm hai vế phương trình đường tròn $x^2 + y^2 = r^2$ theo biến $x$:
\begin{equation}
2x + 2y \frac{dy}{dx} = 0 \implies \frac{dy}{dx} = -\frac{x}{y}
\end{equation}
Độ dốc tiếp tuyến của cung tròn có độ lớn là:
\begin{equation}
|m| = \left|\frac{dy}{dx}\right| = \frac{x}{y}
\end{equation}

Trên cung bát phân khảo sát, hoành độ $x$ bắt đầu từ $0$ và tăng dần đến khi $x = y$, do đó ta luôn có bất đẳng thức:
\begin{equation}
0 \le x \le y \implies 0 \le \frac{x}{y} \le 1 \implies |m| \le 1
\end{equation}

Vì độ lớn độ dốc $|m| \le 1$, tốc độ biến thiên của tung độ $y$ luôn chậm hơn hoặc bằng tốc độ biến thiên của hoành độ $x$. Nếu ta chọn $y$ làm biến tăng đều, bước nhảy của $x$ sẽ lớn hơn 1 pixel ($\Delta x > 1$), khiến các pixel trên đường tròn bị ngắt quãng, tạo ra các lỗ hổng thị giác.

Do đó, \textbf{bắt buộc $x$ phải là biến tăng đều từng bước nguyên đơn vị}:
\begin{equation}
x_{k+1} = x_k + 1
\end{equation}
Khi $x$ tăng đúng 1 pixel, tung độ thực tế $y$ giảm một khoảng nhỏ hơn hoặc bằng 1 pixel ($0 \le -\Delta y \le 1$). Vì vậy, pixel kế tiếp chỉ có đúng hai khả năng lựa chọn: giữ nguyên tung độ ($y_{k+1} = y_k$) hoặc giảm đúng 1 đơn vị ($y_{k+1} = y_k - 1$). Điều này đảm bảo đường cong tròn hoàn toàn liên tục và kết dính trong lân cận 8 pixel.

\subsection{Chứng minh Hình học và Đại số của Thuật toán Điểm Giữa (Midpoint Circle Algorithm)}

Thuật toán Điểm Giữa (Midpoint Circle Algorithm) do Jack E. Bresenham phát triển~\citep{bresenham1977linear} giải quyết bài toán chọn pixel tối ưu bằng cách đánh giá vị trí hình học của điểm chính giữa hai pixel ứng viên:

Giả sử tại bước $k$, thuật toán đã chọn pixel $(x_k, y_k)$. Tại bước $k+1$, hoành độ là $x_{k+1} = x_k + 1$. Hai pixel ứng viên là:
\begin{align}
\text{Pixel trên (hướng Đông, East):} &\quad E(x_k + 1, y_k) \\
\text{Pixel dưới (hướng Đông Nam, SouthEast):} &\quad SE(x_k + 1, y_k - 1)
\end{align}

Điểm chính giữa (Midpoint) $M$ nằm giữa hai pixel ứng viên $E$ và $SE$ có tọa độ:
\begin{equation}
M = \left( x_k + 1, y_k - \frac{1}{2} \right)
\end{equation}

Tham số quyết định tại bước $k$ chính là giá trị của hàm đường tròn tính tại điểm giữa $M$:
\begin{equation}
p_k = f_{\text{circle}}(M) = f_{\text{circle}}\left( x_k + 1, y_k - \frac{1}{2} \right) = (x_k + 1)^2 + \left( y_k - \frac{1}{2} \right)^2 - r^2
\end{equation}

Khoảng cách từ gốc tọa độ $O(0, 0)$ đến điểm giữa $M$ là $d_M = \sqrt{(x_k + 1)^2 + (y_k - 1/2)^2}$, do đó:
\begin{equation}
p_k = d_M^2 - r^2
\end{equation}

\textbf{Chứng minh hai trường hợp quyết định:}
\begin{itemize}
    \item \textbf{Trường hợp 1: $p_k < 0 \iff d_M^2 < r^2 \iff d_M < r$:}
    Điểm giữa $M$ nằm ở phía trong đường tròn thực tế. Điều này đồng nghĩa với việc đường tròn thực tế đi qua phía trên điểm giữa $M$. Khoảng cách từ tâm đến pixel trên $E$ là $d_E > d_M$, trong khi khoảng cách đến pixel dưới là $d_{SE} < d_M$. Vì bán kính đường tròn $r$ thỏa mãn $d_M < r < d_E$, nên cung tròn nằm kẹp giữa $M$ và $E$. Do đó, khoảng cách từ cung tròn đến pixel trên $E$ nhỏ hơn khoảng cách đến pixel dưới $SE$:
    \begin{equation}
    |d_E - r| < |d_{SE} - r|
    \end{equation}
    Như vậy, pixel trên $E(x_k + 1, y_k)$ nằm gần đường tròn hơn, và ta chọn $y_{k+1} = y_k$.
    
    \item \textbf{Trường hợp 2: $p_k \ge 0 \iff d_M^2 \ge r^2 \iff d_M \ge r$:}
    Điểm giữa $M$ nằm ở phía ngoài hoặc nằm chính xác trên đường biên đường tròn. Điều này đồng nghĩa với việc đường tròn thực tế đi qua phía dưới điểm giữa $M$. Khi đó, bán kính đường tròn $r$ thỏa mãn $d_{SE} < r \le d_M$, tức là cung tròn nằm kẹp giữa $SE$ và $M$. Do đó, khoảng cách từ cung tròn đến pixel dưới $SE$ nhỏ hơn hoặc bằng khoảng cách đến pixel trên $E$:
    \begin{equation}
    |d_{SE} - r| \le |d_E - r|
    \end{equation}
    Như vậy, pixel dưới $SE(x_k + 1, y_k - 1)$ nằm gần đường tròn hơn (hoặc cách đều), và ta chọn $y_{k+1} = y_k - 1$.
\end{itemize}

\textbf{Công thức truy hồi cập nhật $p_{k+1}$:}
Tại bước kế tiếp $k+1$, tham số quyết định được tính tại $x_{k+1} + 1$:
\begin{equation}
p_{k+1} = f_{\text{circle}}\left( x_{k+1} + 1, y_{k+1} - \frac{1}{2} \right) = (x_{k+1} + 1)^2 + \left( y_{k+1} - \frac{1}{2} \right)^2 - r^2
\end{equation}
Lấy hiệu số $p_{k+1} - p_k$:
\begin{align}
p_{k+1} - p_k &= \left[ (x_{k+1} + 1)^2 - (x_k + 1)^2 \right] + \left[ \left( y_{k+1} - \frac{1}{2} \right)^2 - \left( y_k - \frac{1}{2} \right)^2 \right] \nonumber \\
&= (x_{k+1} + 1)^2 - x_{k+1}^2 + (y_{k+1}^2 - y_k^2) - (y_{k+1} - y_k) \nonumber \\
&= 2x_{k+1} + 1 + (y_{k+1}^2 - y_k^2) - (y_{k+1} - y_k)
\end{align}

Từ đây, ta thu được hai công thức cập nhật tương ứng:
\begin{itemize}
    \item Nếu $p_k < 0$: Chọn $y_{k+1} = y_k$, biểu thức trở thành:
    \begin{equation}
    p_{k+1} = p_k + 2x_{k+1} + 1
    \end{equation}
    \item Nếu $p_k \ge 0$: Chọn $y_{k+1} = y_k - 1$, biểu thức trở thành:
    \begin{equation}
    p_{k+1} = p_k + 2x_{k+1} + 1 - 2y_{k+1}
    \end{equation}
\end{itemize}

\subsection{Chứng minh Thay thế Khử Số thực: \texorpdfstring{$p_0 = \frac{5}{4} - r$}{p0 = 5/4 - r} bằng \texorpdfstring{$p_0' = 1 - r$}{p0' = 1 - r} khi \texorpdfstring{$r$}{r} Nguyên}

Tại bước khởi tạo ban đầu, thuật toán bắt đầu từ đỉnh trên cùng của đường tròn:
\begin{equation}
(x_0, y_0) = (0, r)
\end{equation}
Giá trị tham số quyết định ban đầu $p_0$ được tính tại điểm giữa đầu tiên $M_0 = (1, r - 1/2)$:
\begin{equation}
p_0 = f_{\text{circle}}\left(1, r - \frac{1}{2}\right) = 1^2 + \left(r - \frac{1}{2}\right)^2 - r^2 = 1 + r^2 - r + \frac{1}{4} - r^2 = \frac{5}{4} - r
\end{equation}

Biểu thức này chứa phân số $\frac{1}{4}$, khiến phép tính khởi tạo ban đầu đòi hỏi số thực:
\begin{equation}
p_0 = 1 - r + \frac{1}{4}
\end{equation}

Để chuyển toàn bộ thuật toán về số học nguyên, ta đặt biến phụ:
\begin{equation}
p_k' = p_k - \frac{1}{4}
\end{equation}
Khi đó, tại bước khởi tạo $k = 0$:
\begin{equation}
p_0' = p_0 - \frac{1}{4} = \left( \frac{5}{4} - r \right) - \frac{1}{4} = 1 - r
\end{equation}
Vì bán kính $r$ là một số nguyên dương ($r \in \mathbb{Z}^+$), nên $p_0' = 1 - r$ chắc chắn là một số nguyên thuần túy.

Hơn thế nữa, xét độ biến thiên giữa hai bước lặp liên tiếp $\Delta p = p_{k+1} - p_k$:
\begin{equation}
p_{k+1}' - p_k' = \left( p_{k+1} - \frac{1}{4} \right) - \left( p_k - \frac{1}{4} \right) = p_{k+1} - p_k
\end{equation}
Độ biến thiên này tại mọi bước đều là số nguyên ($2x_{k+1} + 1$ hoặc $2x_{k+1} + 1 - 2y_{k+1}$). Do đó, $p_k'$ luôn luôn là số nguyên với mọi bước lặp $k \ge 0$.

Bây giờ ta tiến hành chứng minh rằng việc thay thế $p_k$ bằng $p_k'$ không làm thay đổi bất kỳ quyết định rẽ nhánh nào:

Ta có quan hệ:
\begin{equation}
p_k = p_k' + \frac{1}{4} \quad \text{với } p_k' \in \mathbb{Z}
\end{equation}

\textbf{1. Chứng minh chiều điều kiện $p_k \ge 0 \iff p_k' \ge 0$:}
\begin{itemize}
    \item Giả sử $p_k' \ge 0$. Vì $p_k'$ là số nguyên, các giá trị khả dĩ của $p_k'$ là $\{0, 1, 2, 3, \dots\}$. Khi đó:
    \begin{equation}
    p_k = p_k' + \frac{1}{4} \ge 0 + \frac{1}{4} = \frac{1}{4} > 0 \implies p_k \ge 0
    \end{equation}
    \item Ngược lại, giả sử $p_k' < 0$. Vì $p_k'$ là số nguyên, giá trị lớn nhất mà $p_k'$ có thể đạt được là $-1$. Khi đó:
    \begin{equation}
    p_k = p_k' + \frac{1}{4} \le -1 + \frac{1}{4} = -\frac{3}{4} < 0 \implies p_k < 0
    \end{equation}
    Theo nguyên lý phản chứng, nếu $p_k \ge 0$ thì bắt buộc $p_k'$ không thể nhỏ hơn $0$, tức là $p_k' \ge 0$.
\end{itemize}

\textbf{2. Chứng minh chiều điều kiện $p_k < 0 \iff p_k' < 0$:}
Tương tự, từ chứng minh trên ta suy ra ngay lập tức:
\begin{equation}
p_k < 0 \iff p_k' < 0
\end{equation}

\textbf{Kết luận:}
Dấu của $p_k$ và dấu của $p_k'$ hoàn toàn đồng nhất tại mọi bước lặp của thuật toán:
\begin{equation}
p_k \ge 0 \iff p_k' \ge 0 \quad \text{và} \quad p_k < 0 \iff p_k' < 0
\end{equation}
Do đó, ta hoàn toàn có thể thay thế khởi tạo $p_0 = \frac{5}{4} - r$ bằng $p_0' = 1 - r$. Toàn bộ thuật toán Midpoint Circle được thực thi $100\%$ bằng các phép toán số nguyên, khử sạch hoàn toàn sự hiện diện của số thực trên phần cứng.

\newpage
\section{Tuần 5: Mô hình Hình học Camera \& Khôi phục Không gian 3D}
\label{sec:week5}

\subsection{Tổng quan nội dung Tuần 5}
Phần này dành cho các mô hình hình học máy ảnh và tái tạo không gian 3 chiều.

\begin{itemize}
    \item Mô hình Pinhole Camera và các tham số nội/ngoại (Intrinsic \& Extrinsic Parameters).
    \item Hiệu chuẩn Camera (Camera Calibration).
    \item Thị giác hai mắt (Stereo Vision) và Epipolar Geometry.
    \item Structure from Motion (SfM) sơ bộ.
\end{itemize}

\subsection{Câu hỏi \& Bài tập}
\textit{(Nội dung câu hỏi và bài giải chi tiết sẽ được cập nhật tại đây)}

\newpage
\section{Tuần 6: Kiến trúc Mạng Nơ-ron Tích chập (CNN) \& Phân loại Ảnh}
\label{sec:week6}

\subsection{Tổng quan nội dung Tuần 6}
Phần này dành cho mô hình mạng tích chập sâu và bài toán phân loại hình ảnh (Image Classification).

\begin{itemize}
    \item Các thành phần cơ bản của CNN: Convolution, Pooling, Activation (ReLU), Fully-Connected Layers.
    \item Các kiến trúc kinh điển: AlexNet, VGG, ResNet, EfficientNet.
    \item Kỹ thuật huấn luyện: Data Augmentation, Transfer Learning, Regularization.
\end{itemize}

\subsection{Câu hỏi \& Bài tập}
\textit{(Nội dung câu hỏi và bài giải chi tiết sẽ được cập nhật tại đây)}

\newpage
\section{Tuần 7: Phát hiện Đối tượng (Object Detection)}
\label{sec:week7}

\subsection{Tổng quan nội dung Tuần 7}
Phần này dành cho các thuật toán phát hiện đối tượng trong ảnh.

\begin{itemize}
    \item Phương pháp Hai giai đoạn (Two-stage Detectors): R-CNN, Fast R-CNN, Faster R-CNN.
    \item Phương pháp Một giai đoạn (One-stage Detectors): YOLO family (YOLOv1 - YOLOv8), SSD, RetinaNet.
    \item Các chỉ số đánh giá: IoU, Precision-Recall Curve, mAP (Mean Average Precision).
\end{itemize}

\subsection{Câu hỏi \& Bài tập}
\textit{(Nội dung câu hỏi và bài giải chi tiết sẽ được cập nhật tại đây)}

\newpage
\section{Tuần 8: Phân đoạn Ngữ nghĩa \& Phân đoạn Thực thể}
\label{sec:week8}

\subsection{Tổng quan nội dung Tuần 8}
Phần này dành cho bài toán Semantic Segmentation, Instance Segmentation và Panoptic Segmentation.

\begin{itemize}
    \item Semantic Segmentation: Fully Convolutional Networks (FCN), U-Net, DeepLab series.
    \item Instance Segmentation: Mask R-CNN.
    \item Panoptic Segmentation \& Foundation Models: Segment Anything Model (SAM).
\end{itemize}

\subsection{Câu hỏi \& Bài tập}
\textit{(Nội dung câu hỏi và bài giải chi tiết sẽ được cập nhật tại đây)}

\newpage
\section{Tuần 9: Theo dõi Đối tượng \& Phân tích Video (Visual Tracking)}
\label{sec:week9}

\subsection{Tổng quan nội dung Tuần 9}
Phần này dành cho bài toán phân tích chuỗi khung hình video và theo dõi đối tượng.

\begin{itemize}
    \item Single Object Tracking (SOT): Kalman Filter, Mean-Shift, Correlation Filters (KCF), SiamRPN.
    \item Multiple Object Tracking (MOT): Tracking-by-Detection, DeepSORT, ByteTrack.
    \item Nhận dạng hành động (Action Recognition) từ video: 3D CNNs, Two-Stream Networks.
\end{itemize}

\subsection{Câu hỏi \& Bài tập}
\textit{(Nội dung câu hỏi và bài giải chi tiết sẽ được cập nhật tại đây)}

\newpage
\section{Tuần 10: Truy vấn Hình ảnh \& Học Biểu diễn Đặc trưng (Image Retrieval)}
\label{sec:week10}

\subsection{Tổng quan nội dung Tuần 10}
Phần này dành cho bài toán tìm kiếm hình ảnh dựa trên nội dung (CBIR - Content-Based Image Retrieval) và Metric Learning.

\begin{itemize}
    \item Metric Learning: Contrastive Loss, Triplet Loss, ArcFace.
    \item Đánh chỉ mục và tìm kiếm láng giềng gần nhất: Approximate Nearest Neighbors (ANN), FAISS.
    \item Ứng dụng trong nhận diện khuôn mặt và nhận diện người qua camera (Person Re-Identification).
\end{itemize}

\subsection{Câu hỏi \& Bài tập}
\textit{(Nội dung câu hỏi và bài giải chi tiết sẽ được cập nhật tại đây)}

\newpage
\section{Tuần 11: Mô hình Sinh Thị giác (Generative Models: GAN, VAE, Diffusion)}
\label{sec:week11}

\subsection{Tổng quan nội dung Tuần 11}
Phần này dành cho các mô hình sinh hình ảnh và video hiện đại (Generative AI for Computer Vision).

\begin{itemize}
    \item Autoencoders (AE) \& Variational Autoencoders (VAE).
    \item Mạng đối nghịch tạo sinh (Generative Adversarial Networks - GANs): DCGAN, Pix2Pix, CycleGAN, StyleGAN.
    \item Mô hình khuếch tán (Diffusion Models): DDPM, DDIM, Latent Diffusion, Stable Diffusion.
    \item Sinh hình ảnh từ câu lệnh văn bản (Text-to-Image Synthesis).
\end{itemize}

\subsection{Câu hỏi \& Bài tập}
\textit{(Nội dung câu hỏi và bài giải chi tiết sẽ được cập nhật tại đây)}

\newpage
\section{Tuần 12: Đa phương thức (Multimodal AI), 3D Vision \& Tổng kết BTLN}
\label{sec:week12}

\subsection{Tổng quan nội dung Tuần 12}
Phần này dành cho các mô hình đa phương thức tiên tiến, đồ họa 3D hiện đại và tổng kết báo cáo bài tập lớn.

\begin{itemize}
    \item Vision-Language Models (VLM): CLIP, BLIP, LLaVA, GPT-4V.
    \item 3D Vision \& Rendering hiện đại: Neural Radiance Fields (NeRF), 3D Gaussian Splatting (3DGS).
    \item Tổng kết, đánh giá và định hướng phát triển trong Khoa học Dữ liệu Thị giác.
\end{itemize}

\subsection{Câu hỏi \& Bài tập}
\textit{(Nội dung câu hỏi và bài giải chi tiết sẽ được cập nhật tại đây)}


\bibliography{BTLN}
\bibliographystyle{BTLN}

\end{document}
